%%%%%%%%%%%%%%%%%%%%%%%%%%%%%%%%%%%%%%%%%%%%%%%%%%%%%%%%%%%%%%%%%%%%%%%%%%%%
%% Author template for Management Science (mnsc) for articles with no e-companion (EC)
%% Mirko Janc, Ph.D., INFORMS, mirko.janc@informs.org
%% ver. 0.95, December 2010
%%%%%%%%%%%%%%%%%%%%%%%%%%%%%%%%%%%%%%%%%%%%%%%%%%%%%%%%%%%%%%%%%%%%%%%%%%%%
\documentclass[mnsc]{informs3}
%\documentclass[mnsc,blindrev]{informs3}
%\documentclass[mnsc,nonblindrev]{informs3} % current default for manuscript submission

\OneAndAHalfSpacedXI
%%\OneAndAHalfSpacedXII % Current default line spacing
%%\DoubleSpacedXII
%%\DoubleSpacedXI

% If hyperref is used, dvi-to-ps driver of choice must be declared as
%   an additional option to the \documentclass. For example
%\documentclass[dvips,mnsc]{informs3}      % if dvips is used
%\documentclass[dvipsone,mnsc]{informs3}   % if dvipsone is used, etc.

% Private macros here (check that there is no clash with the style)
%%%%%%%% Linkage
\usepackage{xurl}
\usepackage{hyperref}
\hypersetup{colorlinks=true,citecolor=blue,urlcolor=blue}

%%%%%%%% Colored underline
\usepackage{ulem}
\usepackage{soul}
\makeatletter
    \newcommand{\coloruline}[2]{%
        \newcommand\temp@reduline{\bgroup\markoverwith
            {\textcolor{#1}{\rule[-0.5ex]{2pt}{0.4pt}}}\ULon}%
        \temp@reduline{#2}%
    }
    \newcommand{\coloruwave}[2]{%
        \UL@protected\def\temp@uwave{\leavevmode \bgroup
        \ifdim \ULdepth=\maxdimen \ULdepth 3.5\p@
        \else \advance\ULdepth2\p@
        \fi \markoverwith{\textcolor{#1}{\lower\ULdepth\hbox{\sixly \char58}}}\ULon}
        \font\sixly=lasy6 % does not re-load if already loaded, so no memory drain.
        \temp@uwave{#2}%
    }
\makeatother

%%%%%%%% Algorithm
\usepackage{algorithm}
\usepackage{algorithmic}

%%%%%%%% Table, Figure and Diagram
%%%%%%%%%% Table
\usepackage{makecell}
\usepackage{tabularx}
\usepackage{longtable}
\usepackage{array}
\usepackage{booktabs}
\usepackage{multirow}
%%%%%%%%%% Figure
\usepackage{float}
\usepackage{graphicx}
\usepackage{graphics}
\usepackage{caption}
\usepackage{subcaption}
%%%%%%%%%% Diagram Figure
\usepackage{tikz}
\usetikzlibrary{arrows.meta, positioning}
\usepackage{varwidth}
%%%%%%%%%% Comments
\usepackage{comment}


% Natbib setup for author-year style
\usepackage{natbib}
 \bibpunct[, ]{(}{)}{,}{a}{}{,}%
 \def\bibfont{\small}%
 \def\bibsep{\smallskipamount}%
 \def\bibhang{24pt}%
 \def\newblock{\ }%
 \def\BIBand{and}%

%% Setup of theorem styles. Outcomment only one.
%% Preferred default is the first option.
\TheoremsNumberedThrough     % Preferred (Theorem 1, Lemma 1, Theorem 2)
%\TheoremsNumberedByChapter  % (Theorem 1.1, Lema 1.1, Theorem 1.2)
\ECRepeatTheorems

%% Setup of the equation numbering system. Outcomment only one.
%% Preferred default is the first option.
\EquationsNumberedThrough    % Default: (1), (2), ...
%\EquationsNumberedBySection % (1.1), (1.2), ...

% For new submissions, leave this number blank.
% For revisions, input the manuscript number assigned by the on-line
% system along with a suffix ".Rx" where x is the revision number.
\MANUSCRIPTNO{}


%%%%%%%%%%%%%%%%
\begin{document}
%%%%%%%%%%%%%%%%

% Outcomment only when entries are known. Otherwise leave as is and
%   default values will be used.
%\setcounter{page}{1}
%\VOLUME{00}%
%\NO{0}%
%\MONTH{Xxxxx}% (month or a similar seasonal id)
%\YEAR{0000}% e.g., 2005
%\FIRSTPAGE{000}%
%\LASTPAGE{000}%
%\SHORTYEAR{00}% shortened year (two-digit)
%\ISSUE{0000} %
%\LONGFIRSTPAGE{0001} %
%\DOI{10.1287/xxxx.0000.0000}%

% Author's names for the running heads
% Sample depending on the number of authors;
% \RUNAUTHOR{Jones}
% \RUNAUTHOR{Jones and Wilson}
% \RUNAUTHOR{Jones, Miller, and Wilson}
% \RUNAUTHOR{Jones et al.} % for four or more authors
% Enter authors following the given pattern:
\RUNAUTHOR{Zhuang et al.}

% Title or shortened title suitable for running heads. Sample:
% \RUNTITLE{Bundling Information Goods of Decreasing Value}
% Enter the (shortened) title:
\RUNTITLE{Bayesian Structural Inference for Dynamic Crowdsourcing Contests}
\RRHFirstLine{}
\LRHFirstLine{}
\RRHSecondLine{\bf\theRUNAUTHOR:\enskip {\it\theRUNTITLE}}
\LRHSecondLine{\bf\theRUNAUTHOR:\enskip {\it\theRUNTITLE}}

% Full title. Sample:
% \TITLE{Bundling Information Goods of Decreasing Value}
% Enter the full title:
\TITLE{Bayesian Structural Inference for Dynamic Crowdsourcing Contests}

% Block of authors and their affiliations starts here:
% NOTE: Authors with same affiliation, if the order of authors allows,
%   should be entered in ONE field, separated by a comma.
%   \EMAIL field can be repeated if more than one author
\ARTICLEAUTHORS{%
\AUTHOR{Linsheng Zhuang}
\AFF{Institute of Operations Research and Analytics\\
	National University of Singapore, Singapore\\
	Center for Intelligent Computing\\
	Great Bay University, China\\
	\EMAIL{zhuangls@gbu.edu.cn}}
\AUTHOR{Jussi Keppo}
\AFF{NUS Business School \& Institute of Operations Research and Analytics\\
	National University of Singapore, Singapore\\
	\EMAIL{keppo@nus.edu.sg}} %, \URL{}}
\AUTHOR{Zhi Chen}
\AFF{NUS Business School\\
	National University of Singapore, Singapore\\
	\EMAIL{bizczhi@nus.edu.sg}}
% Enter all authors
} % end of the block

\ABSTRACT{%
Online crowdsourcing contests generate rich dynamic behavior, but their core strategic state is latent: contestants react to noisy public leaderboard feedback rather than true relative performance.
We develop a tractable dynamic game with Bayesian filtering and a corresponding structural estimation framework to recover the latent competition state, effort dynamics, participant effort-cost heterogeneity, innovation volatility, and signal precision.
The model admits a closed-form Markov Perfect Equilibrium, enabling Bayesian estimation in a high-frequency dynamic setting.
For validation, we use a path-wise holdout design: private-leaderboard path values are excluded from the estimation likelihood and subsequently used for diagnostic tests, although final private ranks are used to define the focal rivalry.
Contestants do not observe these private-leaderboard path values during the contest, but they are generated from the same submissions and scoring rule as the public leaderboard.
Using data from 73 Kaggle contests, we find that contests with higher estimated innovation risk have more volatile private-leaderboard paths.
For signal precision, the evidence is weaker but directionally consistent: higher estimated precision is associated with smaller public-private leaderboard disagreement.
These results provide diagnostic evidence on the estimated primitives and illustrate how the estimated model can be used to study counterfactual changes in prize levels and contest duration.
% Enter your abstract
}%

% Sample
%\KEYWORDS{deterministic inventory theory; infinite linear programming duality;
%  existence of optimal policies; semi-Markov decision process; cyclic schedule}

% Fill in data. If unknown, outcomment the field
\KEYWORDS{Bayesian Statistics, Structural Estimation, Crowdsourcing Contest, Dynamic Contest}
%\HISTORY{......}

\makeatletter
\def\theARTICLETOP{}
\makeatother

\maketitle
%%%%%%%%%%%%%%%%%%%%%%%%%%%%%%%%%%%%%%%%%%%%%%%%%%%%%%%%%%%%%%%%%%%%%%

% Samples of sectioning (and labeling) in MNSC
% NOTE: (1) \section and \subsection do NOT end with a period
%       (2) \subsubsection and lower need end punctuation
%       (3) capitalization is as shown (title style).
%
%\section{Introduction.}\label{intro} %%1.
%\subsection{Duality and the Classical EOQ Problem.}\label{class-EOQ} %% 1.1.
%\subsection{Outline.}\label{outline1} %% 1.2.
%\subsubsection{Cyclic Schedules for the General Deterministic SMDP.}
%  \label{cyclic-schedules} %% 1.2.1
%\section{Problem Description.}\label{problemdescription} %% 2.

% Text of your paper here

\newpage
\section{Introduction}

% Crowdsourcing contests
Crowdsourcing contests are widely used by firms to source solutions to complex problems from external contributors.
% platforms
Platforms such as Kaggle\footnote{Kaggle: \url{https://www.kaggle.com}}, Topcoder\footnote{Topcoder: \url{https://www.topcoder.com}}, DrivenData\footnote{DrivenData: \url{https://www.drivendata.org}}, and AIcrowd\footnote{AIcrowd: \url{https://www.aicrowd.com}} host thousands of these contests, attracting data scientists, engineers, and developers to compete by submitting solutions to real-world problems.
% how kaggle runs
On Kaggle, for instance, a typical contest spans several weeks, during which participants compete for monetary prizes by repeatedly submitting solutions, observing their relative ranking on a public leaderboard, and adjusting their strategies in response to evolving feedback and the observed actions of competitors.

% Contest structure
The central empirical challenge in these online crowdsourcing contests is that the true competitive state is unobserved during the contest.
Participants observe noisy public leaderboard updates rather than true relative performance, and they choose effort based on filtered beliefs rather than on the latent state itself.
The econometrician faces an analogous partial-observability problem during estimation, making it difficult to recover strategic behavior and incentives from contest-phase data alone.

Motivated by this identification problem, we develop a dynamic contest model and structural Bayesian estimation framework tailored to settings such as Kaggle, where contestants choose continuous effort while learning from noisy public leaderboard signals. 
Applying the framework to the Meta-Kaggle dataset \citep{megan_risdal_timo_bozsolik_2022}, we recover latent competition states and structural primitives from contest-phase public information and evaluate them using post-contest private-leaderboard paths excluded from the likelihood.


\begin{figure}[tbp]
	\center
	\noindent\includegraphics[scale=1.2]{kaggledata_diagram_simplified.pdf}
	\caption{Structure of a Typical Kaggle Contest}
	\label{kaggledata-diagram-simplified}
	\begin{minipage}{\textwidth}
{\footnotesize
\medskip
\textit{Note:}
This figure illustrates the structure of a typical Kaggle contest and the information flow available to contestants.
Rectangles and rounded rectangles represent observed data and latent states, respectively.
The private leaderboard is visible only to the contest organizer during the contest and is revealed to all participants only after it ends.
During the contest period, contestants can observe only the submission times of each player and the public leaderboard.
}
\end{minipage}
\end{figure}


% Framework
To illustrate the framework in greater detail, Figure~\ref{kaggledata-diagram-simplified} depicts the key components and their interactions in a typical Kaggle contest.
% leaderboard
One notable feature is the existence of two leaderboards: a \textit{private} leaderboard, which is hidden until the contest ends and ultimately determines the final ranking and prize allocation, and a \textit{public} leaderboard, which is visible to all participants during the contest.\footnote{
This design serves two purposes.
First, the public leaderboard provides participants with real-time feedback.
Second, by keeping the private leaderboard concealed, the platform discourages overfitting to the public metric and preserves uncertainty, thereby sustaining competitive tension throughout the contest.
}
% submission
Leaderboard displays are contestant-specific and update asynchronously.
% bayesian update
After observing the public leaderboard, contestants update their beliefs and form a perceived state of the competition.
% effort
Based on these beliefs, contestants choose effort according to the model’s Markov perfect equilibrium (MPE) strategies.
% submission
These effort levels shape submission intensity. 
Each submission generates a new score for the submitting contestant, while the displayed score of the opponent remains at its most recent submitted value.
% uncertainty
The private leaderboard, in particular, evolves according to the real-time effort level and an innovation shock.


% Structural estimation
We estimate the structural parameters of the model using a Bayesian approach, based on contestants’ submission times and the observed \textit{public} leaderboard trajectories from the Meta-Kaggle database.
% data
The estimation relies solely on information available to participants during the contest and leverages the equilibrium structure implied by the dynamic model.
% posteriors
This procedure recovers latent competition states, effort dynamics, participant cost parameters, innovation uncertainty, and the precision of the public measurement process. Under the maintained steady-state approximation, signal precision primarily governs filtering and the observation system rather than equilibrium-effort comparative statics.


% Diagnostic validation
To assess whether the estimated model captures meaningful contest primitives, we use private-leaderboard paths to evaluate two predictions illustrated in Figure~\ref{kaggledata-diagram-simplified}: 
higher estimated innovation risk should be associated with greater private-path volatility, 
whereas higher estimated signal precision should be associated with smaller public-private leaderboard gaps. 
Although private-leaderboard path values are excluded from estimation, final private ranks are used to select the focal pair. We therefore interpret the validation exercises as path-wise holdout diagnostics conditional on ex post pair selection.
To strengthen the credibility of these diagnostics, we also conduct time-split exercises that relate primitives estimated from early public contest data to private-path outcomes in the later contest period, thereby reducing the concern that the full-path relationships merely reflect common path-level variation.


% Results
Empirically, we carefully select 73 high-quality contests from the Meta-Kaggle database and preprocess the data to ensure consistency across estimation methods.
Conditional on the focal-pair selection rule, we estimate contest-specific innovation risk and signal precision using only public contest-phase information.
The validation results show that structurally estimated innovation risk is positively associated with private-leaderboard volatility, including in a time-split diagnostic.
For signal precision, the evidence is weaker but directionally consistent: higher estimated signal precision is associated with smaller public-private leaderboard disagreement, and the positive relationship between structural and benchmark precision estimates also appears under an alternative focal-pair specification.
Accordingly, our empirical interpretation of $\lambda$ is primarily as a measurement and filtering primitive.


% Contributions
This paper makes two main contributions.
% 1. model
First, we develop a dynamic game-theoretic model and a Bayesian structural framework for contests with noisy public feedback and latent performance dynamics.
Relative to contest models centered on discrete submit-or-stop decisions, our approach models continuous effort accumulation and makes belief updating from noisy signals structurally central.
Building on the equilibrium structure in \citet{Ryvkin2022Fight}, we add noisy public signals and Bayesian filtering while retaining a closed-form equilibrium for estimation.
% 2. estimation
Second, we use private-leaderboard path values, which are excluded from the likelihood, to construct post-estimation validation targets. Because final private ranks enter focal-pair selection, we interpret these exercises as path-wise holdout diagnostics conditional on ex post pair selection.
This approach allows us to assess whether the estimated structure captures innovation uncertainty and signal precision rather than merely fitting observed contest-phase behavior.
The results provide post-estimation diagnostic evidence on the recovered contest primitives and highlight the value of the model for mechanism evaluation and platform design.
Taken together, these contributions advance our understanding of online competitive dynamics and provide a foundation for future empirical work on digital contest design.



\subsection{Related Literature}

This study is most closely related to the following strands of literature.

\textbf{Contest Design.}
%% Static Contest Design
% definition
A rich literature on \textit{static} contest design \citep{ales2017optimal_award} investigates how to structure competitive environments to elicit performance, typically through a contest success function mapping effort to winning probabilities \citep{skaperdas1996contest}.\footnote{There are exceptions, such as beauty or barrier contests.}
Because our setting is dynamic with repeated feedback, this static branch serves mainly as background for the mechanism-design questions we build on, including prize allocation, participation rules, and information disclosure \citep{moldovanu2001optimal, ales2021number_of_players, bergemann2022optimal}.


%% Dynamic Contest Design
% Transition
While these studies offer valuable insights into incentive design, they largely overlook how participants dynamically adjust their efforts in response to evolving competition.
% Dynamic Contest
In contrast, this paper focuses on \textit{dynamic} contests, where strategic decisions unfold over time.
Two main modeling approaches have been used to capture such dynamic strategic interactions:

% One Way: Poisson Arrival of Innovation Success
One strand of the literature models dynamic contests as multistage processes, where participants exert effort over time while facing uncertainty about intermediate progress.
This structure facilitates the analysis of the mid-contest policy instruments available to the contest designer, such as feedback disclosure and screening mechanisms.
Specifically, in a typical two-stage setting \citep{aoyagi2010information, mihm2019feedback}, the initial stage captures either feasibility uncertainty or foundational experimentation, with progress revealed through milestone feedback provided by the contest organizer.
Participants exert effort in both stages, and the final output reflects the cumulative impact of early performance and subsequent improvement.
Similarly, \citet{bimpikis2019designing} study a two-stage model in which agents exert continuous effort to complete Stage A (potentially infeasible) and Stage B (feasible), with breakthroughs governed by Poisson processes.
In all these models, the contest designer may strategically disclose interim feedback.
Another example is \citet{Khorasani2023screening}, who analyze a two-stage contest in which solvers invest exploratory effort to generate viable ideas, with only those passing an imperfect screening mechanism proceeding to execution.

% Another Way of Modelling
The second approach, more closely related to our work, treats output as a stochastic process whose drift depends on the participant’s effort.
This line of research originates from the tug-of-war framework, first formalized by \citet{Harris1987Race} as a one-dimensional simplification of a multi-stage R\&D race.
Subsequent studies build on this foundation by modeling each player’s progress as a Brownian motion.
\citet{budd1993Duopoly}, for example, analyze a duopoly innovation race where the state evolves as a Brownian motion driven by the difference in effort, and characterize an approximate equilibrium.
\citet{Moscarini2007Contest} advance this approach by modeling the game state as the output gap between players and derive closed-form equilibrium strategies under pure-strategy play.
\citet{Ryvkin2022Fight} adapt their framework to a dynamic contest with a fixed deadline, providing a closed-form solution for the Markov perfect equilibrium.

{\color{blue}
\citet{huang_2024_designing} study the training–test split in data-science contests as a design instrument: intermediate splits make ranking statistically more reliable, but also give contestants the strongest incentive to submit high-variance, suboptimal models. Using recent Kaggle contests, they find evidence consistent with this incentive effect, so that the organizer may accurately select among inaccurate solutions.
}


%% Online Crowdsourcing
% this paper
This paper studies online crowdsourcing contests, which typically feature fixed deadlines, repeated submissions, and public real-time feedback, all governed by platform-specific rules and evaluation metrics.
% our foundation
Hence, our analysis builds on the model by \citet{Ryvkin2022Fight}, whose simplifications yield a closed-form equilibrium that makes structural estimation in a dynamic setting analytically and computationally tractable.
We depart from their full-information setting by introducing a noisy public signal and Bayesian filtering, thereby bringing partial observability into a tractable dynamic contest model.



\textbf{Structural Estimation of Contests.}
Beyond theoretical modeling, a growing body of research employs structural estimation techniques to recover key parameters of contest environments using observed data.

% Static contests
In the structural estimation of static contests, when a contest success function (CSF) is explicitly specified, the primary goal is to estimate its functional form or key parameters \citep{hwang2012technology, kang2016policy, huang2021structural}.
In contrast, for contests without a well-defined or smooth CSF, such as beauty contests, researchers typically rely on equilibrium conditions, most notably the first-order condition.
For example, \citet{yoganarasimhan2016estimation} proposes a two-step semiparametric estimator, resembling the EM algorithm, to recover both unobserved auction heterogeneity and bidder cost distributions.
Similarly, in auction-like models \citep{guerre2000optimal, shakhgildyan2022nonparametric}, nonparametric estimation of the underlying valuation structure often serves as a first step, providing the foundation for subsequent structural estimation.

% Dynamic contests
Most dynamic structural estimation exercises rely on the Markov Perfect Equilibrium (MPE) of a dynamic game.
However, when the model is calibrated to closely fit real-world data, the MPE typically lacks a closed-form solution.
For example, \citet{jiang2022feedback} develop a finite-horizon dynamic game with discrete time, state, and action spaces to model creators’ strategic behavior in response to feedback during online logo design contests.
Using data from 810 contests, they implement a two-step maximum likelihood estimation (MLE) procedure under the assumption of ex-ante homogeneous contests and players: the first step estimates action-specific state transition probabilities, while the second recovers utility parameters.

Specific to online crowdsourcing environments,
\citet{zhang2019structural} build a dynamic structural model to study how participants in crowdsourcing contests learn and choose contests over time.
Their model shows that competing against superstars accelerates learning by providing more informative feedback.
Using Bayesian updating and a hierarchical model, they estimate how individuals trade off monetary rewards, learning opportunities, and participation costs, based on data from Topcoder.

Most closely related to our work, \citet{lemus2021dynamic} structurally estimate contest behavior using Kaggle data.
They propose a finite-horizon dynamic game in which, after each submission, players decide whether to continue or exit, essentially solving an optimal stopping problem.
Players can only make this decision immediately after completing a submission, which arrives at random intervals and incurs a privately drawn cost.
Their estimation proceeds in two stages: the first uses maximum likelihood and the EM algorithm to estimate primitive parameters such as the distribution of player types and scores; the second applies a conditional choice probability (CCP)-based estimator to identify the cost distribution.
Our model differs in two key respects:
(1) we treat submissions as the outcome of accumulated effort rather than as discrete costly actions; and
(2) we explicitly model belief updating from noisy signals as a structural channel that shapes both behavior and identification, whereas their CCP-based approach relies on a coarser behavioral approximation.







The structure of the paper is as follows:
Section~\ref{contest-theory} introduces our theoretical model of online crowdsourcing contests.
Section~\ref{sect-bayes-framework} presents the Bayesian estimation framework.
In Section~\ref{sect-synthetic}, we describe the data-generating process and apply the Bayesian framework to synthetic data.
Section~\ref{sec-kaggle-application} applies the framework to real-world Kaggle data and evaluates the estimated primitives using private-leaderboard path values excluded from the likelihood.
A replication package (code and implementation details) will be provided with the paper.



\section{The Model}\label{contest-theory}

% Effort, Output and Gap
Two players, $i$ and $j$, compete for a prize $\theta$ in a contest.
The winner receives the prize and the loser receives zero.
The contest starts at time zero.
At each time $t\ge0$, player $i$ chooses an effort level $q_{i,t}$ and incurs quadratic cost $C_i(q_{i,t}) = c_i q_{i,t}^2/2$, where lower $c_i$ corresponds to higher ability.
Let $y_t$ denote the \textit{output gap} between players $i$ and $j$ at time $t$, evolving as
\begin{equation}\label{eq-state-dynamics}
	dy_t = (q_{i,t}-q_{j,t})dt + \sigma dW_t
\end{equation}
where $W_t$ is a Brownian motion and $\sigma>0$ measures the innovation risk.
When $\sigma$ is large and the relationship between effort and output is highly uncertain, luck becomes the dominant factor.
Conversely, when $\sigma$ is small and output closely reflects effort, the importance of effort is amplified.

The actual contest is equipped with a submission system in which contestants submit solutions asynchronously over time. Each submission generates new public and private scores for the submitting contestant, while the opponent's most recently submitted scores are carried forward when constructing the corresponding leaderboard gaps.
% Signal
For tractability, we approximate this irregular sequence of submission-triggered updates by a continuous-time public signal of the latent output gap $y_t$.
The precision of this signal is governed by $\lambda$.
The signal evolves according to
\begin{equation}\label{signal}
	dZ_{t} = y_{t}dt + \frac{dB_{t}}{\sqrt{\lambda}}
\end{equation}
where $B_{t}$ is a standard Brownian motion independent of $W_t$, and $\lambda$ controls signal precision.
A larger $\lambda$ implies less noise and more informative feedback; in the empirical mapping, this corresponds to a more informative public leaderboard signal about latent relative performance.
This continuous-time signal is an idealized analogue of the stream of public leaderboard updates observed after submissions.

% Bayesian Player
The information set of both players at time $t \ge 0$ is $I_t \equiv \{Z_s : 0 \le s \le t\}$.
Player $i$ estimates the unknown output gap $y_t$ based on $I_t$.
Let $\tilde{y}_t \equiv E(y_t \mid I_t)$ denote the filtered mean output gap, and let
$S_t \equiv E[(\tilde{y}_t-y_t)^2 \mid I_t]$ denote the conditional mean-squared filtering error.
We normalize the initial perceived gap to $\tilde y_0=0$ throughout the theoretical analysis.
Under standard linear-Gaussian assumptions, the Kalman-Bucy filter (Chapter 1.2 of \citealt{Bensoussan1992Control}; see also \citealt{frogerais2012various,barrau2017invariant}) implies
\begin{align}
	d\tilde{y}_{t} &= (q_{i,t}-q_{j,t})dt + \lambda S_{t}(dZ_{t}-\tilde{y}_{t}dt)
	\label{filtered-x}\\
	\frac{dS_{t}}{dt} &= \sigma^2 - \lambda S_{t}^2
	\label{filtered-S}
\end{align}
Hence, the conditional distribution satisfies $y_t \mid I_t \sim \mathcal{N}(\tilde{y}_t,S_t)$, so the posterior state is fully characterized by $(\tilde y_t,S_t)$.
The steady-state value associated with equation~\eqref{filtered-S} is $\bar S \equiv \sigma/\sqrt{\lambda}$.
For analytical tractability, we set $S_0=\bar S$, so that $S_t \equiv \bar S$ for all $t$.
Fixing \(S_t\) at its steady-state value leaves \(\tilde y_t\) as the only performance state and makes a closed-form equilibrium possible. 
However, the model can no longer capture how filtering uncertainty changes over the course of a contest or how such changes affect effort.
The general solution of equation~\eqref{filtered-S} is reported in Appendix~\ref{app-S-equ}.
Figure~\ref{fig-S-evol} illustrates convergence of $S_t$ to $\bar S$ from arbitrary initial conditions.

\begin{figure}[tbp]
	\centering
	\includegraphics[scale=0.6]{figure_S_evolve.pdf}
	\caption{The evolution of $S_t$ in time $t$, given that $\lambda = 1$ and $\sigma = 1$.} \label{fig-S-evol}
\begin{minipage}{\textwidth}
{\footnotesize
\medskip
\textit{Note:} The plot shows the conditional variance $S_t$ of the filtering error (i.e., the uncertainty about the true performance gap).
It converges quickly from any initial value to its steady state $\bar{S} = \sigma / \sqrt{\lambda}$.
The figure illustrates the convergence of $S_t$ to its steady state.
}
\end{minipage}
\end{figure}



% Contest Deadline
Following \cite{Ryvkin2022Fight}, we consider a dynamic contest with a fixed deadline.
Suppose the contest is terminated when time $t=T>0$.
Under the steady-state restriction $S_t \equiv \bar S$, the payoff-relevant state is fully characterized by $(\tilde y_t,t)$.
Let $\tilde y_T$ denote player $i$'s perceived terminal performance advantage over player $j$.
Although the actual prize is determined by the final private-leaderboard outcome, for tractability we assume that contestants act as if the sign of the filtered mean determines the winner, rather than accounting for the full posterior uncertainty about the private outcome.
Accordingly, player \(i\)'s terminal utility is represented by \(\theta\mathbf 1_{\{\tilde y_T>0\}}\), while player \(j\)'s is represented symmetrically.
This approximation preserves the tractable tug-of-war structure.
At any time $0\le t<T$, player $i$ optimizes her effort level $q_{i,\tau}$ in the remaining contest period $\tau\in[t, T)$ according to the following optimization problem,
\begin{equation}\label{v-def}
	V^i(\tilde{y}_{t}, t ; q_{j,t},\Theta_i) = \max_{\{q_{i,\tau}\}^T_{\tau=t}}
\mathbb{E}\left( \theta\cdot\mathbf{1}_{\tilde y_T>0} - \int^T_tC_i(q_{i,\tau})d\tau \bigg|I_t\right)
\end{equation}
where $\Theta_i \equiv\{\theta, \lambda, \sigma, c_i\}$, subject to constraints (\ref{filtered-x}), (\ref{filtered-S}) and $q_{i,\tau}\ge0$ for all $\tau\in[t,T)$.
The optimization problem for player $j$ is symmetric to that of player $i$; with player labels and cost parameters relabeled, we write $V^j(\tilde{y}_t, t) = V^i(-\tilde{y}_t, t)$.
The corresponding Hamilton-Jacobi-Bellman (HJB) equation for player $i$ is
\begin{equation*}
0 = \max_{q_{i,t}\ge0}\left[-\frac{c_iq_{i,t}^2}{2} + V^i_{y}\cdot\left(q_{i,t}-q_{j,t}\right)+V^i_t + \frac{V^i_{yy}}{2}\lambda \bar{S}^2\right].
\end{equation*}
By definition, we have $\lambda \bar{S}^2 = \sigma^2$.
Assuming an interior solution, substituting the first-order conditions $q_{i,t} = V^i_y/c_i$ and $q_{j,t} = -V^j_y/c_j$ yields the system
\begin{equation*}
\begin{aligned}
\frac{1}{2c_i}(V^i_y)^2 + \frac{1}{c_j}V^i_yV^j_y + V^i_t + V^i_{yy}\frac{\sigma^2}{2} = 0\\
\frac{1}{2c_j}(V^j_y)^2 + \frac{1}{c_i}V^j_yV^i_y + V^j_t + V^j_{yy}\frac{\sigma^2}{2} = 0
\end{aligned}
\end{equation*}
subject to boundary conditions $V^i(-\infty, t) = 0$, $V^i(+\infty, t) = \theta$, $V^j(-\infty, t) = \theta$, $V^j(+\infty, t)=0$, $V^i(\tilde{y}_T, T) = \theta \mathbf 1_{\{\tilde y_T>0\}}$ and $V^j(\tilde{y}_T, T) = \theta \mathbf 1_{\{\tilde y_T<0\}}$.


The Nash equilibrium is summarized in the following theorem.

\begin{theorem}[\citealt{Ryvkin2022Fight}]\label{thm-equilibrium-effort}
Under the terminal-payoff rule above, the game takes the form of a fixed-deadline tug-of-war contest.
In the Markov perfect equilibrium, the players’ efforts $m_{i(j)}(\tilde{y}_t, t) :\mathbb{R}\times[0, T)\to\mathbb{R}_+$ are given by
\begin{equation}\label{eq-equilibrium-effort}
m_{i(j)}(\tilde{y}_t, t) = \frac{e^{-z^2/2}}{\sqrt{2\pi\sigma^2(T-t)}}\cdot \frac{\sigma^2}{2}\left[\gamma(\rho_{i}) + \gamma(\rho_{j})\right]\left[1-\rho(z)^2\right]\left[1 \pm \rho(z)\right]
\end{equation}
where $z = \tilde{y}_t / (\sigma\sqrt{T-t})$, $\rho(z) = \gamma^{-1}\left(\Phi(z)\left[\gamma(\rho_{i})+\gamma(\rho_{j})\right]-\gamma(\rho_{j})\right)$ and
\begin{equation*}
\rho_{i} = \frac{e^{w_{i}}+e^{-w_{j}}-2}{e^{w_{i}}-e^{-w_{j}}},
\quad
\rho_{j} = \frac{e^{w_{j}}+e^{-w_{i}}-2}{e^{w_{j}}-e^{-w_{i}}},
\quad
w_{i(j)} = \frac{\theta}{\sigma^2 c_{i(j)}}.
\end{equation*}
\begin{equation}\label{eq-gamma}
\gamma(u) = \frac{u}{1-u^2} + \frac{1}{2}\ln\frac{1+u}{1-u},\quad u\in(-1,1)
\end{equation}
In the compact notation $m_{i(j)}$, the plus sign in $[1\pm\rho(z)]$ corresponds to player $i$, and the minus sign corresponds to player $j$.
\end{theorem}

In equation~\eqref{eq-equilibrium-effort}, $w_i$ and $w_j$ summarize prize incentives relative to innovation risk, and $\rho_i$ and $\rho_j$ are normalized transformations of these indices.
The equilibrium effort $m_{i(j)}$ can be decomposed into two components: a normal density term $\phi(\tilde{y}_t; 0, \sigma^2(T-t))$ and an amplitude term $K_{i(j)}(z;\theta,\sigma,c_i,c_j)
:= \frac{\sigma^2}{2}\left[\gamma(\rho_i)+\gamma(\rho_j)\right]\left[1-\rho(z)^2\right]\left[1\pm\rho(z)\right]$.
The density term governs how competitive intensity is distributed over the state space, whereas the amplitude term governs the asymmetry of effort across players.

This representation yields three observations useful for the empirical analysis.
First, the parameters $\theta$, $c_i$, and $c_j$ enter the equilibrium effort expression through the amplitude term, so ability and prize incentives primarily affect the level and asymmetry of effort.
Second, $\sigma$ affects both the dispersion of the density term and the amplitude term, so its role is not reducible to a simple scale effect.
Third, as the remaining time $T-t$ shrinks, effort becomes increasingly concentrated near the pivotal state.

The equilibrium characterization is inherited from \citet{Ryvkin2022Fight} after the filtered-state reduction to $(\tilde y_t,t)$ under $S_t\equiv\bar S$.
Additional properties of the equilibrium are discussed in \citet{Ryvkin2022Fight}.
Here, we restrict attention to a brief design discussion under the maintained steady-state specification.

\subsection{Implications for Contest Design}\label{sect-design}

This subsection studies the joint design of the prize $\theta$ and contest duration $T$ when the organizer values the expected best output.
We first characterize total expected effort and show how it contributes to expected best output, and then derive the jointly optimal prize $\theta^*$ and duration $T^*$.

The equilibrium characterization implies the following representation of total expected effort.

\begin{lemma}[\citealt{Ryvkin2022Fight}]\label{lem-total-effort}
The total expected effort at time $t$, defined by $M(\tilde{y}_t, t) := \mathbb{E}\left[\int^T_t m_i(\tilde{y}_s, s) + m_j(\tilde{y}_s, s)ds\right]$, is given by $M(\tilde{y}_t, t) = \sigma \sqrt{T-t} q(z)$, where $q(z)$ is determined by the ordinary differential equation
\begin{equation}
\begin{aligned}
&q''(z) + \left[4A_+(z)e^{-z^2/2}+z\right]q'(z) - q(z) + 4A_-(z)e^{-z^2/2} = 0\\
&A_-(z) = \frac{\gamma(\rho_i) + \gamma(\rho_j)}{2\sqrt{2\pi}}[1 - \rho(z)^2],
\quad A_+(z) = A_-(z)\rho(z)
\end{aligned}
\end{equation}
subject to $q(z)\to0$ as $z\to\pm\infty$.
When $c_i=c_j=c$, so that $w_i=w_j=w$, total expected effort has the leading-order representation
\begin{equation}\label{eq-totaleffort-reduced}
M(\tilde{y}_t, t)
=\sigma\sqrt{T-t}\left[
\frac{2w}{\sqrt{2\pi}}e^{-z^2/2}+o(w)
\right].
\end{equation}
\end{lemma}

Lemma~\ref{lem-total-effort}, established by \citet{Ryvkin2022Fight} and restated here for use in our design analysis, provides an ODE characterization of total expected effort for the general asymmetric case.
For the design analysis below, we use its leading-order symmetric approximation.
This approximation is most informative when $w=\theta/(\sigma^2c)$ is small, so that the prize-cost ratio $\theta/c$ is small relative to innovation variance $\sigma^2$.
Its leading term shows that $w$ governs the overall strength of effort incentives, while the factor $e^{-z^2/2}$ concentrates expected effort in states where the contest remains close.

We next consider the organizer's expected best output.
Define the contestants' individual output processes as
\begin{equation}\label{eq-individual-output-levels}
dx_t^i=m_i(\tilde y_t,t)dt+\frac{\sigma}{\sqrt{2}}dW_t^i,
\qquad
dx_t^j=m_j(\tilde y_t,t)dt+\frac{\sigma}{\sqrt{2}}dW_t^j,
\end{equation}
where $W^i$ and $W^j$ are independent, individual outputs are measured as improvements from the start of the contest so that $x_0^i=x_0^j=0$, and $y_t=x_t^i-x_t^j$.
This representation is consistent with equation~\eqref{eq-state-dynamics}, with $W_t=(W_t^i-W_t^j)/\sqrt{2}$.
Define the organizer's expected best output as
\begin{equation}\label{eq-expected-best-output}
B(\theta,T):=\mathbb E\left[\max\{x_T^i,x_T^j\}\right].
\end{equation}

\begin{lemma}[Expected Best Output]\label{lem-best-output}
The expected best output defined in equation~\eqref{eq-expected-best-output} satisfies the exact decomposition
\begin{equation}\label{eq-best-output-decomposition}
B(\theta,T)
=\frac{1}{2}M(0,0)+\frac{1}{2}\mathbb E|y_T|.
\end{equation}
Moreover, under symmetric costs $c_i=c_j=c$, as $w=\theta/(\sigma^2c)\downarrow0$,
\begin{equation}\label{eq-terminal-gap-leading}
\frac12\mathbb E|y_T|
=\frac{\sigma\sqrt T}{\sqrt{2\pi}}
+O(w^2\sigma\sqrt T).
\end{equation}
\end{lemma}

Equation~\eqref{eq-best-output-decomposition} has a direct interpretation.
The term $M(0,0)$ evaluates the total expected effort in equation~\eqref{eq-totaleffort-reduced} at the initial time $t=0$ under the normalization $\tilde y_0=0$, and therefore measures the expected equilibrium effort accumulated by both contestants over the entire contest period $[0,T]$.
Accordingly, $M(0,0)/2$ is the contribution of aggregate equilibrium effort to expected average terminal output.
The term $\mathbb E|y_T|/2$ is a \textit{selection} component: because the organizer keeps the larger of $x_T^i$ and $x_T^j$, terminal dispersion raises expected best output relative to average output.
Equation~\eqref{eq-terminal-gap-leading} clarifies the leading-order source of this selection component.
When $w=0$, there is no prize-induced effort and the terminal gap is $y_T=\sigma W_T$, where $W_T\sim\mathcal N(0,T)$; hence $\mathbb E|y_T|/2=\sigma\sqrt T/\sqrt{2\pi}$.
Under symmetric costs, an increase in $w$ changes aggregate effort at order $w$, but changes the cumulative effort difference entering $y_T$ only at order $w^2$.
Therefore, the innovation-driven expression remains valid to first order, with the effect of strategic effort differences absorbed by the remainder $O(w^2\sigma\sqrt T)$.
We refer to this leading-order selection gain as the \textit{sampling} component.

Evaluating equation~\eqref{eq-totaleffort-reduced} at the normalized initial state gives $M(0,0)=\bar M(0,0)+o(w\sigma\sqrt T)$, where $\bar M(0,0)=2\theta\sqrt T/(\sigma c\sqrt{2\pi})$ is the leading-order approximation to total expected effort over the entire contest.
Combining Lemmas~\ref{lem-total-effort} and~\ref{lem-best-output}, define the leading-order approximation to expected best output as
\begin{equation}\label{eq-best-output-leading-approximation}
\bar B(\theta,T)
=\frac{\sqrt T}{\sqrt{2\pi}}
\left(\sigma+\frac{\theta}{\sigma c}\right).
\end{equation}
The first term in equation~\eqref{eq-best-output-leading-approximation} is the \textit{sampling} component, and the second term, equal to $\bar M(0,0)/2$, is the first-order \textit{effort} component.

To study contest design, suppose the organizer chooses the prize $\theta\geq0$ and duration $T>0$ to maximize
\begin{equation}\label{eq-designer-best-output-objective}
\Pi(\theta,T)
=V\!\left(\bar B(\theta,T)\right)-\theta-D(T).
\end{equation}
The organizer's value function $V$ is continuous on $[0,\infty)$, twice continuously differentiable on $(0,\infty)$, and bounded above, with $V'>0$, $V''<0$, $V'(0+)=\infty$, and $V'(\infty)=0$.
The prize enters as a direct monetary cost, while $D(T)$ captures the cost of delaying the contest outcome; $D$ is continuously differentiable on $[0,\infty)$, strictly increasing, and convex.

Contest duration $T$ raises both the sampling and effort components of $\bar B(\theta,T)$: a longer contest allows independent innovations to accumulate and gives contestants more time to exert effort.
The optimal duration therefore balances the resulting increase in the organizer's value $V(\bar B)$ against the higher duration cost $D(T)$.
Because equation~\eqref{eq-best-output-leading-approximation} is affine in the prize $\theta$ and $V$ is strictly concave, $\theta$ can be uniquely profiled out of the joint problem.
This yields the following characterization of the jointly optimal prize $\theta^*$ and duration $T^*$.

\begin{proposition}[Joint Design of Prize and Duration]\label{prop-joint-design}
The joint design problem admits at least one optimal pair $(\theta^*,T^*)$ with $0<T^*<\infty$.
Every optimal pair satisfies
\begin{equation}\label{eq-optimal-prize}
\theta^*
=\left[
\sigma c\sqrt{\frac{2\pi}{T^*}}
(V')^{-1}\!\left(\sigma c\sqrt{\frac{2\pi}{T^*}}\right)
-\sigma^2c
\right]_+,
\end{equation}
where $[x]_+=\max\{x,0\}$.
If $\theta^*>0$, the optimal design satisfies
\begin{equation}\label{eq-optimal-duration-condition}
D'(T^*)
=\frac{\theta^*+\sigma^2c}{2T^*}.
\end{equation}
\end{proposition}

\begin{remark}[Leading-Order Design]
Proposition~\ref{prop-joint-design} characterizes design using the leading-order approximation $\bar B$.
The resulting design is internally consistent with the approximation when $w^*=\theta^*/(\sigma^2c)$ is sufficiently small.
\end{remark}

\begin{remark}[Zero-Prize Contests]\label{rem-zero-prize}
When $\theta^*=0$, the symmetric equilibrium, understood by continuous extension as $\theta\downarrow0$, has no prize-induced effort, so $M(0,0)=0$.
Nevertheless, equation~\eqref{eq-best-output-leading-approximation} gives $\bar B(0,T)=\sigma\sqrt T/\sqrt{2\pi}$ and $\Pi(0,T)=V(\sigma\sqrt T/\sqrt{2\pi})-D(T)$.
Thus, a zero-prize contest does not represent the absence of a contest.
Rather, conditional on participation, it represents a contest in which contestants devote no additional effort in response to monetary incentives but may still generate candidate solutions through low-cost, opportunistic experimentation, for example by querying readily available AI tools, and submit the resulting outputs.
The organizer's value then comes entirely from sampling across these baseline random attempts and selecting the best realization.
\end{remark}

\begin{remark}[Comparative Statics]\label{rem-comparative-statics}
The relative curvature $R(B):=-BV''(B)/V'(B)$ measures how rapidly the organizer's marginal value declines as best output increases.
This restriction is not needed for the existence and characterization in Proposition~\ref{prop-joint-design}.
For comparative statics, let $\theta^*(T)$ denote the right-hand side of equation~\eqref{eq-optimal-prize} with $T^*$ replaced by $T$.
If $1/3<R(B)\leq1$ over the relevant output range, then $\theta^*(T)$ is nondecreasing in $T$, and the marginal gross benefit of duration is strictly decreasing.
Together with convexity of $D$, the latter property makes the optimal duration unique.
\end{remark}

Under the terminal-payoff rule and the steady-state restriction $S_t=\bar S$, both performance criteria are neutral with respect to signal precision $\lambda$, because $\lambda\bar S^2=\sigma^2$ and $\lambda$ therefore drops out of the equilibrium effort policy.
This neutrality is specific to the maintained assumptions and need not hold when filtering uncertainty varies over time or contestants respond to posterior winning probabilities.
Nevertheless, $\lambda$ affects belief updating and the noise in the observed public leaderboard, and therefore remains an estimable parameter in the structural model developed below.






\section{Estimation Framework}\label{sect-bayes-framework}

This section presents the data-generating process and the Bayesian estimation procedure.
We first connect the theoretical model to the observable contest data and then describe the likelihood construction and prior specification.

\subsection{Data-Generating Process}

For each Kaggle contest, the observable data from Meta-Kaggle can be grouped into three components.

%% Data: Contest setting
The first component contains contest-level institutional features: contest duration, prize structure, and information disclosure rules.
To match the two-player winner-take-all structure of the model, we restrict the empirical application to contests where a dominant top rivalry is plausible and monetary prize incentives can be mapped to an effective top-two winner-gain structure.

%% Data: Submissions
The second component records submission events of player $i$, denoted by $\{\hat{t}^i_k\}_{k=1}^{N_i}$, where $N_i$ is the total number of submissions.
Conditional on latent effort paths, we model the submissions of players $i$ and $j$ as two inhomogeneous Poisson processes with intensities $\tau_i(t)$ and $\tau_j(t)$. These processes are conditionally independent given the latent effort paths.
For any interval $\mathcal{S}\subset\mathcal{T}$, the expected number of submissions by player $i$ is $\int_{\mathcal{S}}\tau_i(s)ds$.
To link observed submission timing to latent equilibrium effort, we assume for tractability that
\begin{equation}\label{eq-model-intensity}
\tau_i(t) = r \cdot m_i(\tilde{y}_t, t)
\end{equation}
where \(r>0\) converts model effort into submission intensity; the same relationship applies to player \(j\). This proportional relationship is an empirical measurement assumption rather than a result derived from the theoretical game. In estimation, \(r\) is treated as an auxiliary scale parameter rather than a contest primitive, with its range constrained by the prior support specified below.\footnote{
Because $r$ enters the submission-intensity equation rather than the public-signal variance or Brownian innovation covariance directly, its main identification interaction is expected to be with the absolute effort-cost scale. 
A synthetic scale-sensitivity check in the replication package re-estimates the baseline synthetic contest while treating $r$ as unknown. 
The results are consistent with this interpretation, although they do not rule out indirect finite-sample interactions with $\sigma$ or $\lambda$.
}


%% Data: Leaderboard
The third component contains \textit{public} and \textit{private} leaderboard series, denoted by $\hat{x}^{i(j)}_t$ and $x^{i(j)}_t$.
In Kaggle contests, public scores are computed from a released subset, while private scores are computed from withheld data and revealed only after the contest ends.
We define the latent contemporaneous gap as $y_t:=x_t^i-x_t^j$ and the displayed public gap as $\hat y_t:=\hat x_t^i-\hat x_t^j$, where each displayed score equals the contestant's most recently submitted public score. Because submissions are asynchronous, $\hat y_t$ need not compare contemporaneous submissions by the two contestants.
The public gap is the empirical counterpart of the noisy signal in (\ref{signal}), observed through discrete leaderboard updates.
We write this empirical readout as
\begin{equation}\label{eq-model-signal}
dZ_t = \hat y_tdt,
\end{equation}
which should be interpreted as an observation convention linking the public leaderboard gap to the signal path, not as a second structural definition of $Z_t$.
Combining (\ref{signal}) and (\ref{eq-model-signal}) gives the formal white-noise representation
\begin{equation}\label{eq-white-noise-readout}
\hat y_t = y_t+\frac{1}{\sqrt{\lambda}}\xi_t,\qquad \xi_t\equiv \frac{dB_t}{dt},
\end{equation}
where $\xi_t$ denotes Gaussian white noise in the generalized-process sense.
Because white noise is not pointwise observed, equation~(\ref{eq-white-noise-readout}) is not a pointwise sampling equation.
In the empirical data, the public leaderboard is sampled only at submission-triggered update times.
At an empirical update time $t_k$, averaging the white-noise representation over the observed readout interval $h_k=t_k-t_{k-1}$ yields a finite-variance noise term $\varepsilon_k\sim\mathcal{N}(0,1/(\lambda h_k))$.
Approximating the corresponding readout of the latent gap by its endpoint value and substituting the latent state dynamics in \eqref{eq-state-dynamics} gives
\begin{equation}\label{eq-leaderboard-gap}
\hat y_{t_k}
\approx
y_0+\int_0^{t_k}\left[m_i(\tilde{y}_s,s)-m_j(\tilde{y}_s,s)\right]ds
+\sigma W_{t_k}
+\varepsilon_k,
\end{equation}
for update times $t_k$.
Here, $\sigma W_{t_k}$ captures accumulated innovation noise, while $\varepsilon_k$ captures the public-signal noise associated with the leaderboard update.
Thus, larger $\lambda$ implies more precise public feedback.
In the empirical implementation, we normalize $y_0=0$.

% tilde{y}_t
The theoretical model uses a continuous public signal and therefore admits the continuous-time filtering equation in \eqref{filtered-x}. The empirical leaderboard, however, is observed only at asynchronous submission times. An exact empirical implementation would therefore require a hybrid filter with continuous prediction between submissions and discrete measurement updates.

To preserve the tractable one-dimensional state representation used in the equilibrium analysis, we instead employ a continuous-update approximation. Specifically, we carry forward the most recently observed public gap and evolve the perceived gap according to
\begin{equation}\label{eq-filtered-y-update}
d\tilde{y}_{t} = \left[m_i(\tilde{y}_t, t) - m_j(\tilde{y}_t, t)\right]dt + \sqrt{\lambda}\sigma(\hat{y}_t-\tilde{y}_{t}) dt
\end{equation}
with initial condition $\tilde{y}_0 = \mu_0$.
For the empirical implementation only, this initialization relaxes the zero-initial-perceived-gap normalization imposed in the theoretical analysis.
This approximation may overstate learning during long intervals without submissions. The resulting estimates should therefore be interpreted as conditional on the continuous-update approximation rather than as estimates from an exact discrete-observation filter.

We denote the set of estimated parameters by $\Theta := \{c_i, c_j, \sigma, \lambda, \mu_0,r\}$.
Here, $\mu_0$ primarily captures the initial perceived gap and is treated as a nuisance initial-condition parameter rather than a structural design primitive.
The scale parameter $r$ is estimated as an auxiliary measurement parameter that maps latent effort into submission intensity.
Once $\Theta$ is specified, equations (\ref{eq-state-dynamics}), (\ref{eq-leaderboard-gap}), and (\ref{eq-filtered-y-update}) define the working empirical system, while realized trajectories remain stochastic.



\subsection{Modeling Considerations}\label{sect-concerns}

This subsection discusses three practical modeling concerns that are important for interpretation: strategic submission timing, the two-player approximation, and the normality assumption.
The goal is to clarify the scope of the baseline framework and connect each concern to the robustness analyses that follow.

\textbf{Strategic submission timing.}
Beyond the uncertainty caused by partial data disclosure, an additional source of signal noise stems from the timing of leaderboard updates: rankings are only refreshed upon new submissions, rendering the leaderboard uninformative during periods of inactivity.
This lag can bias inference if participants deliberately time or withhold submissions.


Our baseline observation equation assumes submissions are driven by latent effort and does not explicitly model strategic withholding.
Strategic submission timing is an important concern in the literature.
For example, \citet{lemus2021dynamic} adopt a similar simplification by excluding such behavior from their model.
Our estimates should therefore be interpreted for contests where active updating behavior dominates withholding behavior.
As a partial mitigation, we require a minimum number of submissions when selecting representative contestants (Section~\ref{sec-data-clean}), which reduces cases where sparse strategic timing dominates observed dynamics.


\textbf{Selection of Two Players.}
Another concern is the two-player approximation.
Real contests involve many participants, so reducing the environment to two players changes strategic externalities and information structure.
Our interpretation is that the model isolates the dominant top rivalry in contests where competition is concentrated near the frontier.
In Section~\ref{sect-robust-2players}, we conduct a series of robustness checks.
By varying the focal-pair definition, we assess whether the validation results depend on a particular choice of representative rivalry.

\textbf{Gaussian approximation.}
The model relies on Gaussian signal and innovation shocks. Mapping bounded leaderboard scores into the model's state space therefore requires a maintained distributional approximation. The logit transformation is adopted for tractability but does not ensure Gaussianity. Alternative values of its multiplicative scale parameter $a$ assess sensitivity to units, not sensitivity to distributional shape. We therefore interpret the Gaussian specification as a maintained modeling assumption rather than as an empirically validated distributional result.



\subsection{Bayesian Estimation Framework}

Our objective is to develop a Bayesian framework for inferring the unknown parameter set $\Theta := \{ c_i, c_j, \sigma, \lambda, \mu_0,r \}$ using the data introduced above.

Conditional on the selected focal pair, the estimation likelihood uses only contest-phase submission timing and public leaderboard updates.

For numerical integration of the effort and filtering equations, we use a uniform computational grid of length $\Delta_{\mathrm{num}}$. This numerical step size does not determine the observation-noise variance.
By definition, the likelihood function corresponding to the submission times of contestant $i$, $\{t^i_k\}_{k=1}^{N_i}$, is given as follows (a symmetric formulation applies to contestant $j$):
\begin{equation}\label{eq-ihpp-prob}
p\left(\{t^i_k\}_{k=1}^{N_i}\mid\tau_i,\Theta\right)
=
\exp\left\{-\int_{s\in\mathcal T}\tau_i(s)ds\right\}
\prod_{k=1}^{N_i}\tau_i(t^i_k)
\end{equation}
where $\tau_i$ is defined in (\ref{eq-model-intensity}).
The integral term is evaluated on the computational grid.

Next, the public leaderboard gap $\hat y_t$ enters the likelihood only when either contestant submits and receives an update.
We denote these update times by the ordered sequence $(t_1,t_2,\ldots,t_N)$, where $N=N_i+N_j$.
Each $t_k$ is associated with an observed public gap $\hat{y}_{t_k}$.
Let $h_k=t_k-t_{k-1}$, with $t_0=0$, denote the elapsed calendar time associated with the $k$th empirical readout. The quantities $h_k$ are determined by observed submission timing and are distinct from the numerical integration step $\Delta_{\mathrm{num}}$.

To make the conditioning structure explicit, we evaluate the leaderboard likelihood sequentially. Let $\mathcal F_{k-1}$ denote the submission and public-leaderboard history strictly before update time $t_k$. Conditional on $\mathcal F_{k-1}$ and $\Theta$, the empirical filtering recursion generates the one-step-ahead predictive distribution
\[
\hat y_{t_k}\mid\mathcal F_{k-1},\Theta
\sim
\mathcal N\!\left(\mu_{k\mid k-1},V_{k\mid k-1}\right),
\]
Here, $\mu_{k\mid k-1}$ and $V_{k\mid k-1}$ denote the one-step-ahead predictive mean and variance implied by the empirical state and observation equations over $(t_{k-1},t_k]$, conditional only on $\mathcal F_{k-1}$. These moments are computed before incorporating the leaderboard observation at $t_k$.
The leaderboard likelihood is therefore
\begin{equation}\label{eq-obs_gap-prob}
p\left(\{\hat{y}_{t_k}\}_{k=1}^N\mid\Theta\right)
=
\prod_{k=1}^N
\phi\left(\hat y_{t_k}\mid\mu_{k\mid k-1},V_{k\mid k-1}\right).
\end{equation}
The observation at $t_k$ is incorporated into the perceived-state recursion only after its conditional density has been evaluated, so it affects subsequent predicted states and effort paths but not its own one-step-ahead predictive mean.
More generally, the mean would include the latent initial gap $y_0$; in the empirical implementation we normalize $y_0=0$ by measuring subsequent gaps relative to the initial latent-gap origin.
This normalization is distinct from the empirical initialization $\mu_0$, which relaxes the theoretical restriction $\tilde y_0=0$.
Thus, $\mu_0$ enters the likelihood indirectly through the filtering recursion in (\ref{eq-filtered-y-update}), which determines the predictive states, effort paths, and conditional moments used in each factor.
Accordingly, the likelihood uses the public leaderboard values associated with submission updates and preserves the temporal ordering of information.

%% Note: exclude y_t in likelihood
It is worth noting that, although the private leaderboard trajectory $y_t$ is observable ex post in Meta-Kaggle, we deliberately exclude it from the estimation likelihood.
In our framework, agents form beliefs about $y_t$ based solely on $\hat{y}_t$, and their actions are driven by these beliefs.
Conditioning on the realized values of $y_t$ in estimation would effectively bypass the agent’s informational constraints and undermine the role of $\lambda$ in shaping belief formation.
More importantly, it would destroy the path-wise holdout design based on private-leaderboard values.
Treating $y_t$ as latent preserves the model’s behavioral interpretation and enables post-estimation path validation, while final private ranks remain part of focal-pair selection.

%% Approximation
To evaluate (\ref{eq-ihpp-prob}) and (\ref{eq-obs_gap-prob}), we recursively construct $\tilde{y}_t$ and implied efforts $(m_i,m_j)$ from (\ref{eq-equilibrium-effort}) and (\ref{eq-filtered-y-update}).
More precisely, conditional on a parameter vector, perceived states and effort paths are constructed by forward filtering, using each public update only for subsequent states.
However, evaluating $m_{i(j)}(\tilde{y}_t, t)$ via equation~(\ref{eq-equilibrium-effort}) requires numerically approximating the inverse function $\gamma^{-1}$, which poses challenges for the use of gradient-based Markov Chain Monte Carlo (MCMC) sampling methods, such as Hamiltonian Monte Carlo (HMC, \citealt{neal1996bayesian}, \citealt{neal2011mcmc}, \citealt{betancourt2017conceptual}) and the No-U-Turn Sampler (NUTS, \citealt{hoffman2014NUTS}).
Hence, we approximate this inverse function with an analytical form:
\begin{equation}\label{eq-invgamma-approx}
\gamma^{-1}(x) \approx \frac{2}{\pi}\arctan(a \cdot x)
\end{equation}
where $a \approx 0.856$ is obtained by minimizing the infinity norm of the difference between the numerical inverse of $\gamma(\cdot)$ and the approximation $\frac{2}{\pi}\arctan(ax)$.\footnote{Under the $L_1$-norm, the parameter $a = 0.947$; under the $L_2$-norm, the parameter $a=0.895$.}
Figure~\ref{fig-approximation} compares the equilibrium effort function $m_i(\tilde{y}_t, t)$ derived from the approximate analytical form of $\gamma^{-1}$ in equation~(\ref{eq-invgamma-approx}) with that obtained from the numerically accurate solution.
As illustrated, the approximation closely replicates the true function; the approximation error appears negligible relative to posterior uncertainty over the calibration range considered.

\begin{figure}[tbp]
    \centering
    \begin{subfigure}{0.48\textwidth}
        \centering
        \includegraphics[width=\linewidth]{approx_invgamma.pdf}
        \caption{Auxiliary Function $\gamma^{-1}(x)$}
    \end{subfigure}
    \hfill
    \begin{subfigure}{0.48\textwidth}
        \centering
        \includegraphics[width=\linewidth]{approx_equi_effort.pdf}
        \caption{Equilibrium Effort $m_i(\tilde{y}_t, t)$}
    \end{subfigure}
    \caption{Comparison of the Accurate and Approximate Forms}
    \label{fig-approximation}
\begin{minipage}{\textwidth}
{\footnotesize
\medskip
\textit{Note:} Panel (a) compares the exact inverse function $\gamma^{-1}(x)$ of the equilibrium component in (\ref{eq-gamma}) with its analytical approximation $\frac{2}{\pi}\arctan(ax)$ by (\ref{eq-invgamma-approx}), where $a \approx 0.856$ minimizes the infinity-norm approximation error.
Panel (b) compares the resulting equilibrium effort $m_i(\tilde{y}_t, t)$ at $t=0$ (with $T=15$, $\theta=1$, $c_i=c_j=1$, $\sigma=1$) obtained from the analytical approximation versus the exact numerical solution.
}
\end{minipage}
\end{figure}

In the empirical application, we use truncated normal priors with explicit supports:
$c_i,c_j\sim\mathcal{TN}(1,5;[0.5,5])$,
$\sigma\sim\mathcal{TN}(1,5;[0.5,10])$,
$\lambda\sim\mathcal{TN}(1,5;[10^{-6},\infty))$,
$r\sim\mathcal{TN}(10,5;[10^{-2},100])$,
and $\mu_0\sim\mathcal{N}(\hat{y}_0,5^2)$.
Thus, priors are weakly informative for $(c_i,c_j,\sigma,\lambda)$, while $r$ is an auxiliary scale parameter and $\mu_0$ is intentionally more informative to stabilize initialization.


The following lemma formalizes local identification of the structural parameters under the maintained specification, conditional on a fixed effort-to-submission scale $r$.
The parameters are informed by different features of the observed data. Effort asymmetry, submission timing, and early public-gap dynamics provide information about \(c_i\), \(c_j\), and \(\mu_0\). Cross-time covariance in public-gap observations provides information about \(\sigma\), whereas observation-specific signal noise provides information about \(\lambda\).
In the Bayesian implementation above, $r$ is still estimated as an auxiliary measurement scale, but the structural identification statement below is conditional on that scale.
The proof is provided in the Appendix.

\begin{lemma}
\label{lmm-params-identifiability}
For fixed $r$, under the maintained structural specification in (\ref{eq-equilibrium-effort}), (\ref{eq-filtered-y-update}), (\ref{eq-ihpp-prob}) and (\ref{eq-obs_gap-prob}), the parameter vector $(c_i,c_j,\sigma,\lambda,\mu_0)$ is locally identified, provided that the realized public-history support satisfies the full-rank variation condition stated in the proof.
\end{lemma}


\section{Synthetic Experiments}\label{sect-synthetic}

Before applying our estimation procedure to real-world contest data, we evaluate its potential on synthetically generated data.
These synthetic exercises serve three purposes: illustrating the data-generating mechanism, assessing finite-sample parameter recovery, and clarifying which parameters are easier or harder to recover under realistic data conditions.

\begin{algorithm}[!ht]
\caption{Synthetic Data Simulation}
\label{algo-dgp}
\begin{algorithmic}
\STATE \textbf{Input:}
	$\Delta$, $T$, $\theta$, $c_i$, $c_j$, $\sigma$, $\lambda$, $r$, $\tau^\star$, $y_0$, $\mu_0$
\STATE Sample $\{s_k^{i(j)}\}^{N_{i(j)}}_{k=1}$ from Poisson process $(\tau^\star)$ on $[0, T)$
\STATE Sample $\{u_k^{i(j)}\}^{N_{i(j)}}_{k=1}$ from uniform distribution on $[0, 1]$
\STATE Sample series of Brownian motions $(W_t)$ and $(B_t)$
\STATE Initialize $\ell^{i(j)}=0$, $\hat{y}_0 = 0$, $\tilde{y}_t=\mu_0$
\FOR{$t = 0, \Delta, \ldots, T-\Delta$}
    \STATE $m_{i(j)}(\tilde{y}_t, t)$  $\gets$ (\ref{eq-equilibrium-effort}), $\tau^{i(j)}(t)$ $\gets$ (\ref{eq-model-intensity})
    \hfill \COMMENT{Use $\theta$, $c_{i(j)}$, $\sigma$, $r$, $T$}
    \STATE $y_{t+\Delta} \gets$ (\ref{eq-state-dynamics})
    \hfill \COMMENT{Use $m_i-m_j$, $\sigma$, $\Delta$, $W_{t+\Delta}$}
    \STATE $\tilde{y}_{t+\Delta}$ $\gets$ (\ref{eq-filtered-y-update})
    \hfill \COMMENT{Use $m_i-m_j$, $\lambda$, $\sigma$, $\Delta$, and available $\hat{y}_t$}
    \FOR{$s_{k}^{i(j)} \in [t, t+\Delta)$}
        \IF {$u_k^{i(j)} < \tau_{i(j)}(s_{k}^{i(j)}) / \tau^\star_{i(j)}$}
        	    \STATE $\hat{t}^{i(j)}_{\ell} \gets s_{k}^{i(j)}$; $\ell^{i(j)} \gets \ell^{i(j)}+1$
	    \hfill \COMMENT{Accept the submission event}
	    \STATE $\hat{y}_{t+\Delta} \gets $ discretized readout of (\ref{eq-leaderboard-gap})
	    \hfill \COMMENT{Use $\sigma$, $\lambda$, $B_{t+\Delta}$; stored for subsequent filtering}
        \ENDIF
    \ENDFOR
\ENDFOR
\STATE \textbf{Output:} $(y_t)$, $(\tilde{y}_t)$, $(m_{i(j)}(\tilde{y}_t, t))$, $\{\hat{t}_k^{i(j)}\}^{N_{i(j)}}_{k=1}$, $(\hat{y}_{t_k})_{k=1}^{N_i + N_j}$
\end{algorithmic}
\end{algorithm}

A central challenge in generating synthetic data lies in dynamically constructing a point process that conforms to an inhomogeneous Poisson process.
We first generate candidate submission events according to a homogeneous Poisson process over the whole contest duration, using a fixed high intensity $\tau^\star \ge \sup_t \tau_{i(j)}(t)$.
Then, we apply the classical thinning procedure \citep{lewis1979simulation} to determine whether each candidate event is accepted or not.
In plain terms, Algorithm~\ref{algo-dgp} generates candidate submissions, simulates latent and perceived states, accepts submissions via thinning, updates public leaderboard observations, and produces the synthetic sample used in the subsequent Bayesian recovery exercise.

As a baseline example, we consider a synthetic contest with a three-month horizon. The two participants differ slightly in effort costs, with $c_i = 1.2$ and $c_j = 1.5$. The remaining parameters are set to $\theta = 1$, $\sigma = 2.0$, $\lambda = 1.0$, and $r = 15$.
Under this specification, we simulate the latent output gap, the public leaderboard path, the players’ submission times, the perceived output gap, and the corresponding equilibrium effort paths.

\begin{figure}[tbp]
	\noindent\includegraphics[scale=0.75]{synthetic_data.pdf}
	\caption{Synthetic Contest Data over a Three-Month Horizon}
	\label{fig-synthetic_90}
\begin{minipage}{\textwidth}
{\footnotesize
\medskip
\textit{Note:} The blue line represents the latent output gap $y_t$ (private leaderboard), the solid black line represents the public leaderboard $\hat{y}_t$ (updated only upon submissions), and the dashed black line represents the perceived output gap $\tilde{y}_t$ inferred via Bayesian filtering. Vertical ticks mark submission times: ticks above the axis are player $i$'s submissions; ticks below are player $j$'s. Parameters: $\theta = 1.0$, $c_i = 1.2$, $c_j = 1.5$, $\sigma = 2.0$, $\lambda = 1.0$, $r = 15$, $T = 3$ months.
}
\end{minipage}
\end{figure}

Figure~\ref{fig-synthetic_90} shows one realization of the baseline synthetic contest.
The blue line represents the latent output gap $y_t$, which can be interpreted as the \textit{private} leaderboard, while the solid black line represents the \textit{public} leaderboard $\hat y_t$, updated whenever either player submits.
The dashed black line represents the perceived output gap $\tilde y_t$, inferred from the public leaderboard as in equation~\eqref{eq-filtered-y-update}.
The short vertical ticks mark submission times generated by the inhomogeneous Poisson processes.
In this realization, player $i$ submits 28 times and player $j$ submits 18 times.
Between submissions, filtering uncertainty would generally evolve over time; for tractability, however, the synthetic design follows the maintained steady-state specification used throughout the model.

The synthetic recovery exercises use simulation-specific priors, reported in the table notes, and estimate $(c_i,c_j,\sigma,\lambda,\mu_0)$ while fixing the auxiliary effort-to-submission scale $r$ at its simulation value.
When describing a parameter as weakly pinned down in these Bayesian recovery exercises, we use a simple posterior-concentration diagnostic: the posterior standard deviation remains close to its prior standard deviation, and the credible interval shows little shrinkage relative to the prior support.
For $\mu_0$, the prior standard deviation is approximately one, so posterior standard deviations near one indicate limited learning from the realized synthetic data.

Table~\ref{tbl-contest90} presents the Bayesian inference results based on the synthetic data that we display in Figure~\ref{fig-synthetic_90}.
With limited data from a short contest, recovery is uneven across parameters.
In this baseline case, $c_i$, $c_j$, and $\lambda$ are reasonably recoverable, while $\sigma$ is harder to recover and $\mu_0$ remains weakly pinned down (posterior standard deviation close to 1).
This finite-sample pattern is informative and motivates the richer-data exercises below.

\begin{table}[htbp]
\newcolumntype{Y}{>{\centering\arraybackslash}X}
\centering
\caption{Bayesian Estimates from Synthetic Data}\label{tbl-contest90}
\begin{tabularx}{\textwidth}{lYYYYY}
\toprule
\textbf{Parameters} & \textbf{$c_i$} & \textbf{$c_j$} & \textbf{$\sigma$} & \textbf{$\lambda$} & \textbf{$\mu_0$} \\
\addlinespace[0.25ex]
\cline{2-6}
\addlinespace[0.25ex]
(true val.)        & 1.2 & 1.5 & 2.0 & 1.0 & 0.0 \\
\midrule
 Posterior Mean    & 1.089 & 1.751 & 2.801 & 1.105 & 0.005 \\
 Posterior SD       & 0.258 & 0.476 & 0.503 & 0.299 & 0.990 \\
 Post. RMSD            & 0.281 & 0.538 & 0.945 & 0.317 & 0.990 \\
 95\% Interval
		& [0.68, 1.68]
		& [1.01, 2.85]
		& [1.95, 3.93]
		& [0.62, 1.78]
		& [-1.92, 1.94] \\
\bottomrule
\addlinespace[0.5ex]
\end{tabularx}
\begin{minipage}{\textwidth}
{\footnotesize
\textit{Note:}
This table reports the posterior means, posterior standard deviations, posterior root mean squared deviations (Post. RMSD), and 95\% credible intervals from Bayesian estimation using the synthetic contest data shown in Figure~\ref{fig-synthetic_90}. In this single synthetic realization, players $i$ and $j$ submitted 28 and 18 times, respectively. Post. RMSD is defined as $\{(\bar{\theta}-\theta_0)^2+s_\theta^2\}^{1/2}$, where $\theta_0$ is the true simulation value, $\bar{\theta}$ is the posterior mean, and $s_\theta$ is the posterior standard deviation. It is a posterior recovery diagnostic for this realized synthetic sample, not a frequentist repeated-sampling quantity.
The priors are $c_i,c_j\sim\mathcal{TN}(0.5,5;[0.1,5])$, $\sigma\sim\mathcal{TN}(1,5;[0.5,10])$, $\lambda\sim\mathcal{TN}(1,5;[10^{-6},100])$, and $\mu_0\sim\mathcal{TN}(0,1;[-20,20])$; $r$ is fixed at its simulation value.
The 95\% credible interval reports the central posterior interval containing 95\% of the posterior mass.
}
\end{minipage}
\end{table}

As an auxiliary scale-sensitivity check, we re-estimate the same baseline synthetic contest while treating the effort-to-submission scale $r$ as an unknown parameter rather than fixing it at its simulation value.
Table~\ref{tbl-contest90-estimate-r} reports the resulting posterior diagnostic.
The main effect of estimating $r$ is to widen uncertainty about the absolute cost scale: posterior correlations between $r$ and $(c_i,c_j)$ are 0.709 and 0.656, respectively.
By contrast, the corresponding posterior correlations with $\sigma$ and $\lambda$ are 0.055 and 0.009.
Thus, in this synthetic exercise, adding $r$ mainly affects absolute effort-cost scaling, while the innovation-risk and signal-precision primitives remain comparatively stable.

\begin{table}[htbp]
\newcolumntype{Y}{>{\centering\arraybackslash}X}
\centering
\caption{Synthetic Scale-Sensitivity Check with Estimated $r$}\label{tbl-contest90-estimate-r}
\begin{tabularx}{\textwidth}{lYYYYYY}
\toprule
\textbf{Parameters} & \textbf{$c_i$} & \textbf{$c_j$} & \textbf{$\sigma$} & \textbf{$\lambda$} & \textbf{$\mu_0$} & \textbf{$r$} \\
\addlinespace[0.25ex]
\cline{2-7}
\addlinespace[0.25ex]
(true val.)        & 1.2 & 1.5 & 2.0 & 1.0 & 0.0 & 15.0 \\
\midrule
Posterior Mean     & 1.294 & 1.927 & 2.786 & 1.134 & 0.004 & 18.439 \\
Posterior SD       & 0.402 & 0.595 & 0.483 & 0.307 & 0.994 & 4.569 \\
Post. RMSD         & 0.413 & 0.732 & 0.922 & 0.335 & 0.994 & 5.719 \\
95\% Interval
		& [0.64, 2.19]
		& [0.90, 3.17]
		& [1.94, 3.83]
		& [0.65, 1.84]
		& [-1.95, 1.95]
		& [11.09, 28.89] \\
\bottomrule
\addlinespace[0.5ex]
\end{tabularx}
\begin{minipage}{\textwidth}
{\footnotesize
\textit{Note:}
This table repeats the baseline synthetic exercise in Table~\ref{tbl-contest90} but estimates the auxiliary effort-to-submission scale $r$ jointly with $(c_i,c_j,\sigma,\lambda,\mu_0)$.
The prior for $r$ is centered at its simulation value, $r\sim\mathcal{TN}(15,5;[10^{-6},1000])$.
The table is intended as a scale-sensitivity diagnostic for the baseline synthetic realization.
}
\end{minipage}
\end{table}


We next examine how recovery changes with richer data.


\subsection{Parameter Recovery Under Richer Data}

This subsection provides simulation-based recovery diagnostics under two richer-data regimes: pooling replicated contests under fixed structural parameters, and observing longer contests with more submissions.
The goal is practical rather than asymptotic or theorem-based: to assess whether recovery improves in sensible ways as information accumulates.


\textbf{Replications.}
To begin, we investigate recovery in a fixed-design information-scaling exercise.
Starting from the baseline synthetic contest, we replicate the same contest design under identical structural parameters and pool the resulting panels for Bayesian inference.
The submission counts are held fixed at the baseline realization in each replicated panel, so the exercise should not be read as an unconditional simulation of the Poisson count distribution.
Instead, it asks whether posterior recovery improves when the likelihood is given more observations from the same underlying information environment.

The experimental results are presented in Table~\ref{tbl-pool-synthetic-data}.
We consider pools of 10 and 20 replicated panels.
Compared with Table~\ref{tbl-contest90}, posterior RMSD and posterior uncertainty decline for $c_i$, $c_j$, $\sigma$, and $\lambda$, showing that the likelihood contains meaningful identifying content that compounds as more contests are added.
In contrast, $\mu_0$ shows little posterior concentration even with pooled contests.

\begin{table}[htbp]
\newcolumntype{Y}{>{\centering\arraybackslash}X}
\centering
\caption{Bayesian Estimates from Pooled Synthetic Data}\label{tbl-pool-synthetic-data}
\begin{tabularx}{\textwidth}{lYYYYY}
\toprule
\textbf{Parameters} & \textbf{$c_i$} & \textbf{$c_j$} & \textbf{$\sigma$} & \textbf{$\lambda$} & \textbf{$\mu_0$}\\
\addlinespace[0.25ex]
\cline{2-6}
\addlinespace[0.25ex]
(true val.)        & 1.2 & 1.5 & 2.0 & 1.0 & 0.0\\
\midrule
\addlinespace
\multicolumn{6}{l}{\textbf{Pool of 10 Contests}} \\
Posterior Mean  & 1.254 & 1.686 & 1.912 & 1.053 & -0.009\\
Posterior SD     & 0.087 & 0.128 & 0.095 & 0.090 & 0.975\\
Post. RMSD          & 0.102 & 0.225 & 0.130 & 0.104 & 0.975\\
\addlinespace
\multicolumn{6}{l}{\textbf{Pool of 20 Contests}} \\
Posterior Mean  & 1.204 & 1.639 & 1.981 & 0.988 & 0.008\\
Posterior SD     & 0.060 & 0.089 & 0.068 & 0.060 & 1.003\\
Post. RMSD          & 0.060 & 0.165 & 0.071 & 0.062 & 1.003\\
\bottomrule
\addlinespace[0.5ex]
\end{tabularx}
\begin{minipage}{\textwidth}
{\footnotesize
\textit{Note:}
This table reports the posterior means, posterior standard deviations, and posterior root mean squared deviations (Post. RMSD) from Bayesian estimation using \textit{pooled} synthetic panels under identical structural parameters. The pooled exercise fixes the realized submission counts from the baseline synthetic contest in each panel; hence the totals are exactly proportional to the single-contest counts (28 submissions by player $i$ and 18 by player $j$). This design isolates how posterior recovery changes with the amount of information, rather than the additional sampling variation in total Poisson counts.
The priors are $c_i,c_j\sim\mathcal{TN}(1,5;[0.5,5])$, $\sigma\sim\mathcal{TN}(1,5;[0.5,10])$, $\lambda\sim\mathcal{TN}(1,5;[10^{-6},10])$, and $\mu_0\sim\mathcal{TN}(0,1;[-20,20])$; $r$ is fixed at its simulation value.
Post. RMSD is defined as $\{(\bar{\theta}-\theta_0)^2+s_\theta^2\}^{1/2}$ and is computed from the posterior distribution conditional on the realized synthetic sample.
}
\end{minipage}
\end{table}

%%%%%%%%%%%%%%%%%%%%%%%%%%%%%%%%%%%

\textbf{Long-Horizon Contest.}
We then turn to a single-contest regime with a longer time horizon.
Rather than increasing the number of trajectories, we examine how the accumulation of submissions over a longer period improves inference accuracy.

When generating long-horizon contest data, it is important to appropriately lower the effort costs $c_i$ and $c_j$ to ensure sustained effort from the contestants.
This is because effort levels are influenced by the remaining time until the deadline:
according to equation~(\ref{eq-equilibrium-effort}), the further away the deadline, the lower the equilibrium effort.
Hence, this exercise is not a pure horizon comparison: we jointly vary horizon and effective submission intensity to keep the synthetic environment informative.

\begin{table}[htbp]
\newcolumntype{Y}{>{\centering\arraybackslash}X}
\centering
\caption{Bayesian Estimates from Long Term Synthetic Data}\label{tbl-longterm-synthetic-data}
\begin{tabularx}{\textwidth}{lYYYYY}
\toprule
\textbf{Parameters} & \textbf{$c_i$} & \textbf{$c_j$} & \textbf{$\sigma$} & \textbf{$\lambda$} & \textbf{$\mu_0$}\\
\midrule
\multicolumn{6}{l}{\textbf{Contest of 6 Months}} \\
True Value           & 0.6     & 0.75   & 2.0     & 1.0     & 0.0\\
Posterior Mean    & 0.806 & 0.936 & 1.685 & 1.116 & 0.011\\
Posterior SD       & 0.112 & 0.142 & 0.148 & 0.171 & 0.986\\
Post. RMSD            & 0.234 & 0.234 & 0.348 & 0.207 & 0.986\\
\addlinespace
\multicolumn{6}{l}{\textbf{Contest of 12 Months}} \\
True Value           & 0.3     & 0.375 & 2.0     & 1.0     & 0.0\\
Posterior Mean    & 0.333 & 0.436 & 2.067 & 1.101 & 0.012\\
Posterior SD       & 0.027 & 0.046 & 0.097 & 0.138 & 0.994\\
Post. RMSD            & 0.043 & 0.076 & 0.118 & 0.171 & 0.994\\
\bottomrule
\addlinespace[0.5ex]
\end{tabularx}
\begin{minipage}{\textwidth}
{\footnotesize
\textit{Note:}
This table reports the posterior means, posterior standard deviations, and posterior root mean squared deviations (Post. RMSD) from Bayesian estimation using synthetic contest data with longer horizons (6 and 12 months). In the 6-month contest, players $i$ and $j$ submitted 59 and 51 times, respectively; in the 12-month contest, they submitted 100 and 97 times, respectively. To sustain active submission behavior over longer horizons, effort costs $c_i$ and $c_j$ are reduced proportionally (true values: $c_i = 0.6$, $c_j = 0.75$ for 6 months; $c_i = 0.3$, $c_j = 0.375$ for 12 months).
The priors are the same as in the baseline single-contest exercise: $c_i,c_j\sim\mathcal{TN}(0.5,5;[0.1,5])$, $\sigma\sim\mathcal{TN}(1,5;[0.5,10])$, $\lambda\sim\mathcal{TN}(1,5;[10^{-6},100])$, and $\mu_0\sim\mathcal{TN}(0,1;[-20,20])$; $r$ is fixed at its simulation value.
Post. RMSD is defined as $\{(\bar{\theta}-\theta_0)^2+s_\theta^2\}^{1/2}$ and is computed from the posterior distribution conditional on the realized synthetic sample.
}
\end{minipage}
\end{table}

Table~\ref{tbl-longterm-synthetic-data} shows that posterior recovery improves as within-contest information increases, although this improvement reflects both a longer horizon and the lower effort costs used to maintain active submission behavior.

Taken together, the synthetic exercises indicate a consistent recoverability ranking: $c_i$, $c_j$, and $\lambda$ are comparatively well recovered; $\sigma$ is harder in thinner samples but improves with richer data; and $\mu_0$ remains weakly learned according to the posterior-concentration criterion above.




\section{Empirical Application}
\label{sec-kaggle-application}

In this section, we apply the model to real Kaggle contests and assess whether structurally recovered contest primitives align with private-leaderboard-based benchmarks.

We select representative Kaggle competitions from Meta-Kaggle and identify two focal participants $i$ and $j$ in each contest based on submission frequency, activity duration, and final ranking.
Using their submission records and public leaderboard trajectories, we construct the observed data $\{t^i_k\}^{N_i}_{k=1}$, $\{t^j_k\}^{N_j}_{k=1}$ and $\{\hat{y}_{t_k}\}_{k=1}^{N_i+N_j}$ for contest-level Bayesian estimation under a unified prior specification.

The estimation likelihood uses only contest-phase submission timing and public leaderboard trajectories. 
Private-leaderboard path values are excluded from estimation, but final private ranks are used to select the focal pair. 
We therefore interpret the resulting evidence as diagnostic validation conditional on the selected pair, rather than as validation based on a fully independent sample.
Specifically, we compare structural contest-level estimates of $\lambda$ and $\sigma$ against benchmark quantities constructed from private-leaderboard path values excluded from the likelihood.
By the private-leaderboard path, we mean the ex post sequence of private scores associated with the contestants' historical submissions, reconstructed after the contest; it is not a path observed by contestants in real time.
Figure~\ref{kaggledata-diagram} summarizes the data-model structure and validation design.

\begin{figure}[tbp]
	\centering
	\noindent\includegraphics[scale=1.2]{kaggledata_diagram.pdf}
	\caption{Meta-Kaggle Dataset and the Model Structure}
	\label{kaggledata-diagram}
	\begin{minipage}{\textwidth}
{\footnotesize
\medskip
\textit{Note:}
This figure shows the structure of a typical Kaggle contest and the information flow available to contestants.
Rectangles represent observed data, rounded rectangles represent latent states, and ellipses represent contest parameters.
Solid lines depict temporal relationships among time series variables and state variables, while dashed arrows show the influence of contest parameters on the data-generating process.
The private leaderboard $y_t$ is not observed by contestants during the contest. Its path values are excluded from the likelihood and used for post-estimation validation, while final private ranks enter focal-pair selection.
}
\end{minipage}
\end{figure}


\subsection{Preparing Data for Structural Estimation}
\label{sec-data-clean}

To empirically implement the structural estimation framework developed in the previous sections using Kaggle data, the first step is to map the observed data onto the corresponding variables defined by the model.
This subsection details the procedures for selecting competitions, identifying key participants, and incorporating contest-specific information such as prize structures to make the data compatible with the model’s structural assumptions.

% Reward
As a first selection criterion, we focus on monetary prize structure.
We restrict attention to Kaggle competitions offering one, two, or three monetary prizes denominated in U.S. dollars.
Competitions with a single prize naturally conform to the winner-take-all incentive structure.
For competitions with two or three prize tiers, we align the empirical setting with the model by treating the prize difference between the top two contestants as the effective gain from winning.
Because non-winning contestants may still receive symbolic recognition such as medals, focusing on the top-two prize gap helps mitigate confounding incentives and brings the empirical setting closer to the model's winner-take-all benchmark.

% Players
The second criterion identifies two highly active, top-ranked contestants whose participation overlaps substantially during the contest. We select candidate focal contestants beginning from the top of the final private leaderboard. This focuses the two-player model on rivalry among high-performing contestants. Private path values do not enter the estimation likelihood, but final private ranks determine the focal pair; hence, the validation exercises are path-wise holdout diagnostics conditional on ex post pair selection.
To ensure overlapping activity and adequate data for estimation, each contestant must submit at least five times, and the pair must satisfy the temporal-overlap criterion used in sample construction.

While analytically convenient, restricting attention to two participants may limit representativeness, oversimplify the strategic environment, and introduce selection concerns via \textit{ex post} rankings.
We therefore interpret the two-player model as a structural approximation of focal rivalry near the top of contests with concentrated competition, frequent submissions, and overlapping activity windows.

% Leaderboard Display
The third criterion for contest selection concerns the format in which leaderboard scores are presented.
Since Kaggle competitions are based on diverse real-world tasks, the scoring metrics vary accordingly.
Some are based on raw scores (positive real numbers), while others are based on accuracy measures (expressed as percentages or decimals between 0 and 1).
Ignoring these differences would lead to inconsistencies in units across contests, complicating any meaningful cross-contest comparison of parameter estimates.

We restrict attention to contests reporting accuracy-based scores, which provides a common bounded score domain and facilitates a consistent transformation across contests.
The retained focal paths are restricted to score observations satisfying the domain $p_t\in[0.5,1)$ used below.
To map raw leaderboard accuracies into the model's state scale, we first transform each individual leaderboard score using the logit (log-odds) function:
\begin{equation}\label{eq-leaderboard-transform}
x_t = a \cdot \log\left(\frac{p_t}{1-p_t}\right), \quad p_t\in[0.5,1)
\end{equation}
where $p_t$ is an individual leaderboard score, $x_t$ is the corresponding transformed score, and $a$ is a scaling parameter in the transformation.\footnote{
Changing $a$ rescales the transformed score and can affect scale-dependent likelihood and validation quantities, but it does not change scale-invariant shape diagnostics of the transformed distribution.
In this paper, to preserve comparability across contests and avoid extra tuning degrees of freedom, we fix $a=1$ throughout.
}
For the focal pair $(i,j)$, we then construct the model variables as transformed-score gaps: $\hat y_t=\hat x_t^i-\hat x_t^j$ for the public leaderboard and $y_t=x_t^i-x_t^j$ for the private leaderboard.
We construct hourly evaluation grids by carrying forward each contestant's most recently observed transformed score between submissions.
Because submissions are asynchronous, only the submitting contestant’s score is updated, while the opponent’s most recent score is carried forward. 
The observed public gap may therefore compare scores submitted at different times and only approximates the contemporaneous latent performance gap.
Our observation equation treats this gap as an approximate readout of the latent competitive state; this approximation is most credible for focal pairs with frequent submissions and substantial overlap in active periods.
On the domain used here, the log-odds transformation is increasing and convex, so as $p_t$ approaches 1, a given increase in accuracy produces a larger change in the transformed score $x_t$.


\subsection{Private-Leaderboard Primitive Validation}\label{sect-validation}


We select 73 contests from the Meta-Kaggle database, encompassing all competitions concluded before April 15, 2025 that satisfy our eligibility filters on leaderboard availability, public-private score information, and focal-player activity.\footnote{
The Meta-Kaggle database is updated daily, with new competition data added after the contests conclude.
}
For each selected contest, we identify two highly active top-ranked contestants who meet our focal-pair selection criteria.
Appendix~\ref{appendix-statistics} reports the Bayesian estimation results for each contest individually.

\begin{table}[H]
\centering
\caption{Public--Private Leaderboard Agreement}
\label{tab-leaderboard-agreement}
\begin{tabularx}{\textwidth}{>{\raggedright\arraybackslash}Xccc}
\toprule
\textbf{Statistic} & \textbf{Mean} & \textbf{Median} & \textbf{95\% CI for Mean} \\
\midrule
Public--private rank correlation & $0.937$ & $0.984$ & $[0.903,\ 0.963]$ \\
Top-2 overlap rate & $65.1\%$ & $50.0\%$ & $[56.2\%,\ 73.3\%]$ \\
Top-5 overlap rate & $66.6\%$ & $80.0\%$ & $[59.7\%,\ 73.2\%]$ \\
\bottomrule
\end{tabularx}
\par\vspace{1.2ex}
\begin{minipage}{\textwidth}
{\footnotesize
\textit{Note:}
The sample consists of the same 73 Kaggle contests used in the structural analysis. For each contest, we use the official public and private leaderboard ranks in Meta-Kaggle, excluding benchmark teams and teams missing either rank. Public--private rank correlation is the Spearman correlation across all retained teams within a contest. The top-$k$ overlap rate is $|\mathcal T^{pub}_{k}\cap\mathcal T^{pri}_{k}|/k$, where $\mathcal T^{pub}_{k}$ and $\mathcal T^{pri}_{k}$ are the sets of teams ranked in the public and private top $k$, respectively; the ordering of teams within the top-$k$ sets need not coincide. Each statistic is first calculated separately by contest, and the table reports the equally weighted mean and median across the 73 contests. The 95\% confidence intervals pertain to the across-contest means and are percentile intervals obtained from 10,000 nonparametric bootstrap resamples at the contest level.
}
\end{minipage}
\end{table}

We next conduct post-estimation diagnostic validation of the contest-specific primitives using private-leaderboard-based benchmarks.

\subsubsection{Private-Leaderboard Diagnostic for Signal Precision $\lambda$.}
The signal precision parameter $\lambda$ is contest-specific and captures the informativeness of public leaderboard feedback relative to latent performance. We compare the structural estimate obtained from submissions and public leaderboard paths with an approximate benchmark constructed from public-private leaderboard discrepancies.
This benchmark is informative but not fully independent of the maintained structural setup, because it still relies on the transformed leaderboard scale used in the model.

Let $\ell\in\{1,\ldots,L\}$ index contests.
Under the maintained Gaussian discrepancy approximation induced by the discretized signal readout, we construct the following signal-precision benchmark.
Let $\mathcal E_\ell$ denote the ordered set of submission-triggered leaderboard readouts for the focal pair, and let $N^\ell=|\mathcal E_\ell|$.
Let $h_k^\ell=t_k^\ell-t_{k-1}^\ell$ denote the elapsed time associated with readout $t_k^\ell$.
Under the maintained approximation
\[
\hat y_{t_k}^\ell-y_{t_k}^\ell
\sim
\mathcal N\!\left(0,\frac{1}{\lambda_\ell h_k^\ell}\right),
\qquad t_k^\ell\in\mathcal E_\ell,
\]
we construct
\begin{equation*}
\lambda^{\text{bench}}_\ell =
\frac{N^\ell}
{\sum_{t_k^\ell\in\mathcal E_\ell}
h_k^\ell (y_{t_k}^\ell - \hat{y}_{t_k}^\ell)^2}.
\end{equation*}
The benchmark therefore uses only submission-triggered leaderboard readouts rather than repeated hourly carry-forward observations.
\footnote{
As a robustness comparison to the irregular-interval estimator above, we also consider the fixed-grid approximation $\lambda_{\ell,\Delta}^{\mathrm{bench}}=N^\ell/\left\{\Delta\sum_{t_k^\ell\in\mathcal E_\ell}(y_{t_k}^\ell-\hat y_{t_k}^\ell)^2\right\}$. 
Despite this difference, both constructions yield the same qualitative conclusion of a positive but modest cross-contest association; the fixed-grid construction also underlies the alternative focal-pair validation in Section~\ref{sect-robust-2players}.
}
Because asynchronous updates may repeatedly use an opponent's carried-forward score, the discrepancies can be serially dependent; we therefore interpret $\lambda_\ell^{\mathrm{bench}}$ as an approximate signal-precision benchmark.

This benchmark relies on the Gaussian discrepancy approximation applied to the transformed scores in \eqref{eq-leaderboard-transform}.
The transformation is adopted for tractability and should be interpreted as an approximation rather than as a claim of exact normality for every contest. Auxiliary checks with alternative choices of the scale parameter $a$ assess sensitivity to transformed-score units, not sensitivity to distributional shape. The Gaussian discrepancy structure therefore remains a maintained modeling assumption.

We then examine the cross-contest association between the private-leaderboard benchmark $\lambda_\ell^{\mathrm{bench}}$ and the Bayesian structural estimate $\hat\lambda_\ell$.

%%%%%%%%%% Figure & Table %%%%%%%%%%%
\begin{figure}[tbp]
\centering
\begin{minipage}[t]{0.43\textwidth}
	\vspace*{\fill}
	\centering
	\includegraphics[width=\linewidth]{validate_lamb.pdf}
	\vspace*{\fill}
\end{minipage}
\hfill
\begin{minipage}[t]{0.53\textwidth}
	\vspace*{\fill}
	\centering
	\newcolumntype{Y}{>{\centering\arraybackslash}X}
		\begin{tabularx}{\textwidth}{llYY}
			\toprule
			& coef. & stderr & 95\% CI \\
			\midrule
			const.
			& -0.808 & 3.157 & [-7.10, 5.49] \\
			$\log(1 + \hat{\lambda})$
			& 3.273\textbf{**} & 1.081 & [1.12, 5.43] \\
			\midrule
			Obs. & \multicolumn{3}{c}{73} \\
			R$^2$ & \multicolumn{3}{c}{0.114} \\
			Adj. R$^2$ & \multicolumn{3}{c}{0.102} \\
		\bottomrule
		\addlinespace[0.5ex]
	\end{tabularx}
	\begin{minipage}{\textwidth}
	{\footnotesize
	(1) Significance levels: \textbf{***} $p < 0.001$, \textbf{**} $p < 0.01$, \textbf{*} $p < 0.05$.
	(2) OLS residual diagnostics: Omnibus = 0.397, Jarque-Bera = 0.071, Skewness = 0.032, Kurtosis = 3.139.
	(3) Other diagnostic statistics: Condition Number = 27.9, Durbin-Watson = 1.806.
}
	\end{minipage}
	\vspace*{\fill}
\end{minipage}
\caption{$\lambda_\ell^{\text{bench}}$ vs. Bayesian Estimate $\hat\lambda_\ell$}
\label{fig-validate-lambda}
\begin{minipage}{\textwidth}
{\footnotesize
\textbf{Note.} This figure plots the approximate contest-specific signal-precision benchmark $\lambda_\ell^{\mathrm{bench}}$ against the Bayesian structural estimate $\hat{\lambda}_\ell$ for each of the 73 contests. The benchmark is computed from public-private leaderboard discrepancies at submission-triggered readouts as $\lambda_\ell^{\mathrm{bench}} = N^\ell / \sum_{t_k\in\mathcal E_\ell} h_k^\ell (y_{t_k}^\ell - \hat{y}_{t_k}^\ell)^2$, where $N^\ell=|\mathcal E_\ell|$ is the number of readouts used in contest $\ell$ and $h_k^\ell=t_k^\ell-t_{k-1}^\ell$ is the elapsed time associated with readout $t_k^\ell$. Because carried-forward scores may induce dependence across successive readouts, the benchmark should be interpreted as an approximate measure of signal precision. 
The solid line represents the fitted regression line from $\log(1 + \lambda_\ell^{\mathrm{bench}}) = \beta_0 + \beta_1 \log(1 + \hat{\lambda}_\ell) + \varepsilon_\ell$. The positive slope ($\beta_1 = 3.273$, $p < 0.01$) indicates a statistically significant but modest association ($R^2 = 0.114$).
}
\end{minipage}
\end{figure}


Figure~\ref{fig-validate-lambda} indicates a positive association between structural and benchmark estimates of $\lambda$.
The slope on $\log(1+\hat\lambda)$ is positive (3.273, $p<0.01$), but the fit is modest ($R^2\approx0.11$), so this should be interpreted as directional evidence rather than precise contest-level recovery.
This dispersion is consistent with limited sample size and with approximation error in benchmark construction, especially the normality-sensitive transformation step.
Section~\ref{sect-robust-2players} provides a complementary, directional check using an alternative focal-pair definition; the positive pattern for $\lambda$ also appears there.


\subsubsection{Public-Private Leaderboard Disagreement Diagnostics.}
We next provide an additional diagnostic for signal precision based directly on public-private leaderboard disagreement, without first converting the disagreement into a Gaussian precision benchmark.
If $\hat\lambda_\ell$ captures public-signal precision, contests with larger $\hat\lambda_\ell$ should exhibit smaller disagreement between the public leaderboard path and the corresponding private-leaderboard path.
We therefore construct the following contest-level outcome:
\[
D_\ell^{path}
=
\frac{1}{|\mathcal H_\ell|}
\sum_{h\in\mathcal H_\ell}
\left|y_{\ell h}-\hat y_{\ell h}\right|,
\]
where $\mathcal H_\ell$ denotes the set of hourly observation points used in the empirical leaderboard path, and $D_\ell^{path}$ measures the mean absolute public-private gap over the contest path.
The predicted association is therefore negative.

Because the full-path diagnostic uses the same contest path both to estimate $\hat\lambda_\ell$ and to construct $D_\ell^{path}$, it may partly reflect common path-level variation rather than purely predictive content.
We therefore also report a time-split version of this diagnostic.
Let $\hat\lambda_\ell^{early}$ denote the structural estimate obtained using only the first $\tau T_\ell$ hours of contest-phase data, with $\tau=0.75$.
The validation target is then constructed only from the remaining public-private path:
\[
D_\ell^{late}
=
\frac{1}{|\mathcal H_\ell^{late}|}
\sum_{h\in\mathcal H_\ell^{late}}
\left|y_{\ell h}-\hat y_{\ell h}\right|.
\]
Here $\mathcal H_\ell^{late}=\{h\in\mathcal H_\ell:h>\tau T_\ell\}$.
For both diagnostics, we estimate the log-log validation slope linking the discrepancy measure to the corresponding signal-precision estimate.
The full-path specification uses $\log(1+D_\ell^{path})$ and $\log(1+\hat\lambda_\ell)$, while the time-split specification uses $\log(1+D_\ell^{late})$ and $\log(1+\hat\lambda_\ell^{early})$.
The latter asks whether signal precision estimated from early public contest information predicts public-private disagreement measured over the later contest period, conditional on the focal-pair selection rule.

\begin{table}[htbp]
\newcolumntype{Y}{>{\centering\arraybackslash}X}
\centering
\caption{Signal Precision and Public-Private Path Disagreement}
\label{tab-signal-disagreement-validation}
\begin{tabularx}{\textwidth}{lYYYY}
\toprule
\textbf{Specification} & \textbf{Coef.} & \textbf{95\% CI} & \textbf{$R^2$} & \textbf{Spearman $\rho$} \\
\midrule
Full path
& \makecell{$-0.156$\textbf{*}\\($0.063$)} & [$-0.282$, $-0.030$] & $0.079$ & $-0.340$\textbf{**} \\
Time split
& \makecell{$-0.046$\\($0.092$)} & [$-0.225$, $0.134$] & $0.003$ & $-0.220$ \\
\midrule
Obs. & \multicolumn{4}{c}{73} \\
\bottomrule
\addlinespace[0.5ex]
\end{tabularx}
\begin{minipage}{\textwidth}
{\footnotesize
\textit{Note:}
For contest $\ell$, $D_\ell^{path}=|\mathcal H_\ell|^{-1}\sum_{h\in\mathcal H_\ell}|y_{\ell h}-\hat y_{\ell h}|$ measures the mean absolute gap between private and public leaderboard paths.
The full-path row reports the slope from $\log(1+D_\ell^{path})=\alpha+\beta\log(1+\hat\lambda_\ell)+\varepsilon_\ell$, where $\hat\lambda_\ell$ is the baseline structural estimate.
For the time-split row, $\hat\lambda_\ell^{early}$ is estimated using only the first 75\% of contest-phase data, and $D_\ell^{late}=|\mathcal H_\ell^{late}|^{-1}\sum_{h\in\mathcal H_\ell^{late}}|y_{\ell h}-\hat y_{\ell h}|$ is measured using only the remaining 25\% of the private and public leaderboard paths.
The reported coefficient is the slope from $\log(1+D_\ell^{late})=\alpha+\beta\log(1+\hat\lambda_\ell^{early})+\varepsilon_\ell$.
Standard errors are reported in parentheses under the coefficient estimates.
Significance levels: \textbf{***} $p<0.001$, \textbf{**} $p<0.01$, \textbf{*} $p<0.05$.
}
\end{minipage}
\end{table}

Table~\ref{tab-signal-disagreement-validation} reports both specifications.
The full-path diagnostic has the predicted negative slope ($-0.156$, $p<0.05$) and a negative Spearman correlation ($-0.340$, $p<0.01$), consistent with the view that higher estimated signal precision is associated with smaller public-private leaderboard disagreement.
The time-split estimate has the same sign but is smaller and imprecise, so it should be interpreted as directional rather than confirmatory evidence.
Taken together with the benchmark-precision exercise above, these results provide suggestive diagnostic evidence for the $\lambda$ interpretation: the predicted associations appear in the full-sample specifications, but signal precision is not precisely recovered at the contest level.
The weaker time-split result is expected because $\hat\lambda_\ell^{early}$ is estimated from less data and $D_\ell^{late}$ is measured over a shorter validation window, both of which increase sampling variability.
\medskip




\subsubsection{Private-Leaderboard Consistency Benchmark for Innovation Risk $\sigma$.}
Another key contest-specific parameter is $\sigma$, which captures innovation uncertainty and the extent to which contest outcomes are driven by chance rather than effort.
For each contest $\ell$, we estimate $\hat{\sigma}_\ell$ structurally using the two contestants’ public submission records $\{t_k^\ell\}_{k=1}^{N_i^\ell+N_j^\ell}$ and public leaderboard trajectories $\{\hat y_k^\ell\}_{k=1}^{N_i^\ell+N_j^\ell}$.
We then construct a private-leaderboard-based benchmark for $\sigma$ using score increments.

Let $(t^\ell_k,y^\ell_k)_{k=0}^{n_\ell}$ denote the private-leaderboard gap evaluated at the ordered update times used for this benchmark, including the normalized initial observation.
Under the diffusion specification in (\ref{eq-state-dynamics}), the private-leaderboard increment satisfies
$y^\ell_{k} - y^\ell_{k-1} \mid \{m_i^\ell,m_j^\ell\} \sim \mathcal{N}\left(\mu^\ell_k, \sigma^2(t^\ell_k-t^\ell_{k-1})\right)$ for $k=1,\ldots,n_\ell$.
Here, $\mu^\ell_k = \int^{t^\ell_k}_{t^\ell_{k-1}}\left[m^\ell_i(s) - m^\ell_j(s)\right]ds$.
Motivated by the Gaussian increment formula, we compute the following contest-specific consistency benchmark; the corresponding diffusion calculation is provided in Appendix~\ref{appexidx-normal-benchmark}:
\begin{equation*}
\sigma^{\text{bench}}_\ell = \left(\frac{1}{n_\ell}\sum_{k=1}^{n_\ell}\frac{\left[y^\ell_k-y^\ell_{k-1}-\hat\mu^\ell_k\right]^2}{t^\ell_k-t^\ell_{k-1}}\right)^{1/2}
\end{equation*}
The drift adjustment $\hat\mu_k^\ell$ is computed from the Bayesian posterior effort trajectories, so this benchmark is model-assisted rather than independent. In addition, because contestant submissions are asynchronous, consecutive displayed private-gap observations may combine a new score for one contestant with a carried-forward score for the other. The resulting estimator should therefore be interpreted as an approximate consistency benchmark, not as an exact maximum likelihood estimator of latent diffusion volatility. For this reason, we place greater weight on the less model-dependent private-path volatility diagnostics reported below.

As in the analysis above for $\lambda$, we regress the benchmark $\sigma^{\text{bench}}$ on the structural estimate $\hat{\sigma}$ to examine whether a positive relationship exists between the two.

%%%%%%%%%% Figure & Table %%%%%%%%%%%
\begin{figure}[tbp]
\centering
\begin{minipage}[t]{0.43\textwidth}
	\vspace*{\fill}
	\centering
	\includegraphics[width=\linewidth]{validate_sigma.pdf}
	\vspace*{\fill}
\end{minipage}
\hfill
\begin{minipage}[t]{0.53\textwidth}
	\vspace*{\fill}
	\centering
	\newcolumntype{Y}{>{\centering\arraybackslash}X}
	\begin{tabularx}{\textwidth}{llYY}
		\toprule
		& coef. & stderr & 95\% CI \\
		\midrule
		const.
		& 0.054 & 0.116 & [-0.18, 0.29] \\
		$\log(1 + \hat{\sigma})$
		& 0.887\textbf{***} & 0.152 & [0.58, 1.19] \\
		\midrule
		Obs. & \multicolumn{3}{c}{73} \\
		R$^2$ & \multicolumn{3}{c}{0.317} \\
		Adj. R$^2$ & \multicolumn{3}{c}{0.308} \\
		\bottomrule
		\addlinespace[0.5ex]
	\end{tabularx}
	\begin{minipage}{\textwidth}
{\footnotesize
(1) Significance levels: \textbf{***} $p < 0.001$, \textbf{**} $p < 0.01$, \textbf{*} $p < 0.05$.
(2) OLS residual diagnostics: Omnibus = 1.280, Jarque-Bera = 0.774, Skewness = 0.225, Kurtosis = 3.214.
(3) Other diagnostic statistics: Condition Number = 5.08, Durbin-Watson = 1.618.
}
	\end{minipage}
	\vspace*{\fill}
\end{minipage}
\caption{$\sigma^{\text{bench}}$ vs. Bayesian Estimate $\hat\sigma$}
\label{fig-validation-sigma}
\begin{minipage}{\textwidth}
{\footnotesize
\textbf{Note.} This figure plots the private-leaderboard consistency benchmark $\sigma^{\mathrm{bench}}$ against the Bayesian structural estimate $\hat{\sigma}$ for each of the 73 contests. The benchmark is computed from carried-forward private leaderboard increments as $\sigma^{\mathrm{bench}} = \left(\frac{1}{n} \sum_{k=1}^n \frac{(y_k-y_{k-1}-\hat{\mu}_k)^2}{t_k-t_{k-1}}\right)^{1/2}$, where $n$ is the number of adjacent displayed-gap increments and $\hat{\mu}_k$ is computed from posterior effort trajectories. Because asynchronous observations can combine a new score with an opponent's stale score, the quantity is an approximate, model-assisted consistency benchmark rather than an exact diffusion MLE. The dashed line represents the fitted regression line from $\log(1 + \sigma^{\mathrm{bench}}) = \beta_0 + \beta_1 \log(1 + \hat{\sigma}) + \varepsilon$.
}
\end{minipage}
\end{figure}


Figure~\ref{fig-validation-sigma} shows a positive association between the structural estimate and the private-leaderboard volatility benchmark. The relationship is statistically significant, although the benchmark remains approximate because of asynchronous observations and the posterior drift adjustment. We therefore interpret this result as supportive diagnostic evidence for the estimated innovation-risk primitive.

\subsubsection{Private-Leaderboard Volatility Diagnostics.}
The preceding validation of $\sigma$ uses an approximate private-leaderboard consistency benchmark motivated by the normal increment structure of the diffusion model.
That benchmark is parametric and uses the posterior drift adjustment $\hat\mu_k$.
We therefore add a simpler diagnostic based on realized private-leaderboard volatility.
For each contest $\ell$, let $\Delta y_{\ell h}=y_{\ell h}-y_{\ell,h-1}$ denote the increment on the constructed hourly evaluation grid, and define $V_\ell^{path}=\operatorname{sd}\{\Delta y_{\ell h}:h\in\mathcal H_\ell,\ h>1\}$.
This statistic does not require estimating a normal-increment likelihood or subtracting a model-implied drift term.
It asks whether contests with larger estimated innovation risk also exhibit more volatile private-leaderboard paths.
Because the hourly private paths are constructed by carrying scores forward between submissions, the resulting series contains zero increments between updates and discrete jumps at submission times. This volatility measure therefore reflects both changes in submitted scores and the timing of private-leaderboard updates. We interpret it as a less model-dependent path diagnostic, not as a direct measure of latent diffusion volatility.
We estimate this relationship in logs, using $\log(1+V_\ell^{path})$ and $\log(1+\hat\sigma_\ell)$, with the predicted positive coefficient.

We also report a time-split version of the same diagnostic.
Let $\hat\sigma_\ell^{early}$ denote the structural estimate obtained using only the first $\tau T_\ell$ hours of contest-phase data, with $\tau=0.75$.
The validation target is then constructed only from the remaining private-leaderboard path, $V_\ell^{late}=\operatorname{sd}\{\Delta y_{\ell h}:h\in\mathcal H_\ell,\ h>\tau T_\ell\}$.
For the time-split specification, we use $\log(1+V_\ell^{late})$ and $\log(1+\hat\sigma_\ell^{early})$.
Conditional on the focal-pair selection rule, the time-split version provides a stricter temporal separation because the structural estimate uses only early public contest data, while the validation outcome is measured on the later private-leaderboard path.

\begin{table}[htbp]
\newcolumntype{Y}{>{\centering\arraybackslash}X}
\centering
\caption{Private-Leaderboard Volatility Diagnostics}
\label{tab-private-volatility-diagnostics}
\begin{tabularx}{\textwidth}{lYYYY}
\toprule
\textbf{Specification} & \textbf{Coef.} & \textbf{95\% CI} & \textbf{$R^2$} & \textbf{Spearman $\rho$} \\
\midrule
Full path
& \makecell{$0.278$\textbf{***}\\($0.078$)} & [$0.126$, $0.430$] & $0.356$ & $0.492$\textbf{***} \\
Time split
& \makecell{$0.220$\textbf{*}\\($0.101$)} & [$0.022$, $0.418$] & $0.153$ & $0.314$\textbf{**} \\
\midrule
Obs. & \multicolumn{4}{c}{73} \\
\bottomrule
\addlinespace[0.5ex]
\end{tabularx}
\begin{minipage}{\textwidth}
{\footnotesize
\textit{Note:}
For contest $\ell$, $y_{\ell h}$ denotes the private-leaderboard gap on the constructed hourly evaluation grid, and $\Delta y_{\ell h}=y_{\ell h}-y_{\ell,h-1}$ is the corresponding grid increment.
The full-path row reports the slope $\beta$ from $\log(1+V_\ell^{path})=\alpha+\beta\log(1+\hat\sigma_\ell)+\varepsilon_\ell$, where $V_\ell^{path}=\operatorname{sd}\{\Delta y_{\ell h}:h\in\mathcal H_\ell,\ h>1\}$ and $\hat\sigma_\ell$ is the baseline structural estimate.
The time-split row reports the slope $\beta$ from $\log(1+V_\ell^{late})=\alpha+\beta\log(1+\hat\sigma_\ell^{early})+\varepsilon_\ell$, where $\hat\sigma_\ell^{early}$ is estimated using only the first 75\% of contest-phase data and $V_\ell^{late}=\operatorname{sd}\{\Delta y_{\ell h}:h\in\mathcal H_\ell,\ h>0.75T_\ell\}$.
Spearman $\rho$ is the rank correlation between the same dependent variable and structural estimate used in each row.
Standard errors are reported in parentheses under the coefficient estimates.
Significance levels: \textbf{***} $p<0.001$, \textbf{**} $p<0.01$, \textbf{*} $p<0.05$.
}
\end{minipage}
\end{table}

Table~\ref{tab-private-volatility-diagnostics} shows a positive relationship between estimated innovation risk and realized private-leaderboard volatility.
The full-path slope is 0.278 ($p<0.001$), with $R^2=0.356$ and a Spearman correlation of 0.492 ($p<0.001$).
The time-split slope is smaller but remains positive and significant: 0.220 ($p<0.05$), with a Spearman correlation of 0.314 ($p<0.01$).
These results complement the normal-increment consistency benchmark and suggest that the $\sigma$ estimates capture meaningful variation in private-leaderboard volatility.
Because both $\hat\sigma_\ell$ and the private-path volatility measures are expressed in transformed-score units, their cross-contest association may also reflect contest-level scale variation. We therefore interpret these regressions as diagnostics rather than scale-free validation evidence.


\subsection{Counterfactual Best-Output-Preserving Prize-Duration Substitution}

The validation exercises suggest that the estimated primitives contain
meaningful information about the contest environment. We use these primitives to
revisit the design question in Section~\ref{sect-design}. Under the maintained
approximation, prize and duration are the two design variables that directly
affect predicted expected best output through the effort and sampling components.\footnote{
We do not emphasize a baseline information-design counterfactual for
$\lambda$. As shown in Section~\ref{sect-design}, under the maintained
steady-state filtering restriction $S_t\equiv\bar S$, signal precision does not
generate an independent comparative-static effect on predicted expected best output.
Studying optimal information disclosure would require relaxing the steady-state
restriction and allowing filtering uncertainty to evolve endogenously over time.
}
Because the designer's cardinal payoff and the platform cost of longer contests
are difficult to calibrate, we do not solve for a welfare-optimal design.
Instead, we trace the model-implied substitution between prize spending and
contest duration while holding predicted expected best output fixed.

For each contest, let $(\theta^{obs},T^{obs})$ denote the observed prize and
duration, and let
$B_0=B(\theta^{obs},T^{obs})$ denote predicted expected best output under the observed
design. For a duration multiplier $a\geq1$, let $T_a=aT^{obs}$ and define
\[
\theta_a^{\mathrm{BO}}
=
\min\left\{
\theta\in[\theta_{\min},\theta^{obs}]
:
B(\theta,T_a)\geq B_0
\right\},
\]
where $\theta_{\min}=0.01$ thousand USD is the lower prize bound used in the numerical search. For each duration multiplier considered below, direct numerical evaluation confirms that $\theta=\theta^{obs}$ satisfies the best-output constraint, so the feasible set is nonempty.
The required prize change is
$100(\theta_a^{\mathrm{BO}}-\theta^{obs})/\theta^{obs}$, so negative values indicate prize
reductions that keep predicted expected best output at least at the observed-design
baseline. This exercise is not a claim that longer contests are costless or
welfare improving. Rather, it reports the iso-best-output prize-duration frontier:
how much prize spending can be substituted by a small, externally specified increase in duration.

For the counterfactual calculation, we combine the exact decomposition in Lemma~\ref{lem-best-output} with the leading-order selection component in equation~\eqref{eq-terminal-gap-leading} and evaluate predicted expected best output as $B(\theta,T)=M(\tilde y_0,0)/2+\sigma\sqrt T/\sqrt{2\pi}$. At each candidate design $(\theta,T)$, total expected effort $M(\tilde y_0,0)$ is obtained from the general asymmetric ODE in Lemma~\ref{lem-total-effort}, using the contest-specific estimates of $(c_i,c_j,\sigma,\mu_0)$ and setting $\tilde y_0=\mu_0$. Thus, the calculation retains the general asymmetric numerical effort prediction while approximating the selection component at leading order.

Table~\ref{tab-counterfactual-policy} reports the substitution frontier for duration changes close to the observed design. For each positive duration increase considered, the numerical search identifies a lower prize that preserves predicted expected best output in all 73 contests. A 1\% increase in duration allows a median prize reduction of 1.56\%, while a 10\% increase allows a median reduction of 13.56\%. The mean percentage reductions exceed the medians in magnitude, so we emphasize the medians as the more representative summary. At the 10\% duration increase, one contest reaches the numerical prize floor $\theta_{\min}$ and consequently delivers predicted expected best output strictly above the observed-design baseline.


\begin{table}[htbp]
\newcolumntype{Y}{>{\centering\arraybackslash}X}
\centering
\caption{Counterfactual Best-Output-Preserving Prize-Duration Substitution}
\label{tab-counterfactual-policy}
\scriptsize
\begin{tabularx}{\textwidth}{lYYYYYY}
\toprule
\textbf{Duration multiplier} & \textbf{Lower prize} & \textbf{$\geq1\%$ lower} & \textbf{$\geq5\%$ lower} & \textbf{$\geq10\%$ lower} & \textbf{Mean change} & \textbf{Median change} \\
\midrule
$1.00$ & $0/73$ & $0/73$ & $0/73$ & $0/73$ & $0.00\%$ & $0.00\%$ \\
$1.01$ & $73/73$ & $73/73$ & $3/73$ & $2/73$ & $-2.44\%$ & $-1.56\%$ \\
$1.02$ & $73/73$ & $73/73$ & $12/73$ & $3/73$ & $-4.52\%$ & $-3.09\%$ \\
$1.05$ & $73/73$ & $73/73$ & $73/73$ & $19/73$ & $-10.29\%$ & $-7.37\%$ \\
$1.10$ & $73/73$ & $73/73$ & $73/73$ & $65/73$ & $-17.92\%$ & $-13.56\%$ \\
\bottomrule
\addlinespace[0.5ex]
\end{tabularx}
\begin{minipage}{\textwidth}
{\footnotesize
\textit{Note:}
This table reports best-output-preserving prize-duration substitution results for the 73 contests with complete observed design inputs and structural estimates.
For duration multiplier $a$, the counterfactual duration is $T_a=aT^{obs}$.
For each contest, we compute the minimum prize $\theta_a^{\mathrm{BO}}\in[\theta_{\min},\theta^{obs}]$ thousand USD such that $B(\theta_a^{\mathrm{BO}},T_a)\geq B(\theta^{obs},T^{obs})$. At each candidate design, predicted expected best output is evaluated as $B(\theta,T)=M(\tilde y_0,0)/2+\sigma\sqrt T/\sqrt{2\pi}$. Total expected effort is evaluated as $M(\tilde y_0,0)=\sigma\sqrt{T}\,q(\tilde y_0/(\sigma\sqrt{T}))$, where $q$ solves the ordinary differential equation in Lemma~\ref{lem-total-effort}; the sampling term is the leading-order selection component in equation~\eqref{eq-terminal-gap-leading}. We use the contest-specific estimates of $(c_i,c_j,\sigma,\mu_0)$, set $\tilde y_0=\mu_0$, and use the same analytical approximation to $\gamma^{-1}$ as in estimation. The numerical lower bound is $\theta_{\min}=0.01$ thousand USD.
``Lower prize'' counts contests with $\theta_a^{\mathrm{BO}}<\theta^{obs}$.
The next three columns count contests in which the required prize is at least 1\%, 5\%, or 10\% below the observed prize.
Mean and median changes report $100(\theta_a^{\mathrm{BO}}-\theta^{obs})/\theta^{obs}$, so negative values denote required prize reductions.
}
\end{minipage}
\end{table}





\section{Robustness}

This section examines the sensitivity of our main validation results to key modeling choices.
We focus on two concerns highlighted in Section~\ref{sect-concerns}: the selection of the two-player focal pair and the normality assumption underlying the leaderboard transformation. For each robustness check, we re-estimate the model and compare the resulting validation patterns to the baseline findings reported in Section~\ref{sect-validation}.

\subsection{Selection of Two Contestants}\label{sect-robust-2players}

The baseline empirical analysis pairs two highly active top-ranked contestants in each contest.
To assess whether the validation results are sensitive to this focal-pair choice, we re-estimate the model under an alternative pairing rule.

%\textbf{Alternative Pairing: Second vs. Third Final Rank.}
As a robustness check, we estimate the model on the second- and third-ranked contestants in each contest, based on final private leaderboard ranks.
This specification shifts the focal rivalry away from the winner and tests whether the validation pattern is driven by one specific top-pair definition.

%%%%%%%%%% Figure & Table %%%%%%%%%%%
\begin{figure}[tbp]
\centering
\begin{minipage}[t]{0.43\textwidth}
	\vspace*{\fill}
	\centering
	\includegraphics[width=\linewidth]{validate_lamb_robust23.pdf}
	\vspace*{\fill}
\end{minipage}
\hfill
\begin{minipage}[t]{0.53\textwidth}
	\vspace*{\fill}
	\centering
	\newcolumntype{Y}{>{\centering\arraybackslash}X}
	\begin{tabularx}{\textwidth}{llYY}
		\toprule
		& coef. & stderr & 95\% CI \\
		\midrule
		const.
		& -0.290 & 3.819 & [-7.91, 7.32] \\
		$\log(1 + \hat{\lambda})$
		& 3.421\textbf{**} & 1.244 & [0.94, 5.90] \\
		\midrule
		Obs. & \multicolumn{3}{c}{73} \\
		R$^2$ & \multicolumn{3}{c}{0.096} \\
		Adj. R$^2$ & \multicolumn{3}{c}{0.084} \\
		\bottomrule
		\addlinespace[0.5ex]
	\end{tabularx}
	\begin{minipage}{\textwidth}
{\footnotesize
(1) Significance levels: \textbf{***} $p < 0.001$, \textbf{**} $p < 0.01$, \textbf{*} $p < 0.05$.
(2) OLS residual diagnostics: Omnibus = 3.192, Jarque-Bera = 2.424, Skewness = -0.411, Kurtosis = 3.350.
(3) Other diagnostic statistics: Condition Number = 36.9, Durbin-Watson = 2.236.
}
	\end{minipage}
	\vspace*{\fill}
\end{minipage}
\caption{Robustness Check (2nd--3rd Pair): $\lambda_\ell^{\text{bench}}$ vs. Bayesian Estimate $\hat\lambda_\ell$}
\label{fig-validation-lambda-robust23}
\begin{minipage}{\textwidth}
{\footnotesize
\textbf{Note.} This figure reports the validation comparison under an alternative focal pair: the second- and third-ranked contestants (based on final private leaderboard ranks) instead of the top two. The dashed line is the fitted regression $\log(1 + \lambda_\ell^{\mathrm{bench}}) = \beta_0 + \beta_1 \log(1 + \hat{\lambda}_\ell) + \varepsilon_\ell$, with $\beta_1 = 3.421$ ($p < 0.01$) and $R^2 = 0.096$. The positive slope remains significant, though the fit is modest, so the result is best interpreted as directional evidence under an alternative focal-pair definition.
}
\end{minipage}
\end{figure}

For signal precision, the alternative focal-pair specification based on the second- and third-ranked contestants again yields a positive and significant relationship, although the explanatory power remains limited ($R^2=0.096$).
This check asks whether changing the focal-pair definition overturns the qualitative validation pattern.
It does not: the sign and significance are preserved, while the low fit reinforces the more cautious interpretation that signal precision is harder to recover precisely at the contest level.

%%%%%%%%%% Figure & Table %%%%%%%%%%%
\begin{figure}[tbp]
\centering
\begin{minipage}[t]{0.43\textwidth}
	\vspace*{\fill}
	\centering
	\includegraphics[width=\linewidth]{validate_sigma_robust23.pdf}
	\vspace*{\fill}
\end{minipage}
\hfill
\begin{minipage}[t]{0.53\textwidth}
	\vspace*{\fill}
	\centering
	\newcolumntype{Y}{>{\centering\arraybackslash}X}
	\begin{tabularx}{\textwidth}{llYY}
		\toprule
		& coef. & stderr & 95\% CI \\
		\midrule
		const.
		& 0.117 & 0.130 & [-0.14, 0.38] \\
		$\log(1 + \hat{\sigma})$
		& 0.501\textbf{*} & 0.212 & [0.08, 0.92] \\
		\midrule
		Obs. & \multicolumn{3}{c}{73} \\
		R$^2$ & \multicolumn{3}{c}{0.071} \\
		Adj. R$^2$ & \multicolumn{3}{c}{0.059} \\
		\bottomrule
		\addlinespace[0.5ex]
	\end{tabularx}
	\begin{minipage}{\textwidth}
{\footnotesize
(1) Significance levels: \textbf{***} $p < 0.001$, \textbf{**} $p < 0.01$, \textbf{*} $p < 0.05$.
(2) OLS residual diagnostics: Omnibus = 12.72, Jarque-Bera = 13.66, Skewness = 1.02, Kurtosis = 3.50.
(3) Other diagnostic statistics: Condition Number = 7.84, Durbin-Watson = 1.87.
}
	\end{minipage}
	\vspace*{\fill}
\end{minipage}
\caption{Robustness Check (2nd--3rd Pair): $\sigma^{\text{bench}}$ vs. Bayesian Estimate $\hat\sigma$}
\label{fig-validation-sigma-robust23}
\begin{minipage}{\textwidth}
{\footnotesize
\textbf{Note.} This figure replicates the validation exercise in Figure~\ref{fig-validation-sigma} using an alternative focal pair: the second- and third-ranked contestants (based on final private leaderboard ranks) instead of the top two. The dashed line is the fitted regression $\log(1 + \sigma^{\mathrm{bench}}) = \beta_0 + \beta_1 \log(1 + \hat{\sigma}) + \varepsilon$, with $\beta_1 = 0.501$ ($p < 0.05$) and $R^2 = 0.071$. The positive slope remains significant, though the association is weaker than in the baseline ($\beta_1 = 0.887$, $R^2 = 0.317$), indicating that $\sigma$ recovery is directionally robust but sensitive to the choice of focal pair.
}
\end{minipage}
\end{figure}

For innovation risk, the positive and significant association also survives this alternative pairing rule, although the fit is weaker than in the baseline specification.
Together, these results suggest that the signs and significance patterns of the main validation exercises are not driven solely by one focal-pair definition. At the same time, the strength of the association, especially for $\sigma$, is sensitive to the selected pair.


\subsection{Normality Caveat and Interpretation}\label{sect-robust-normality}

The multiplicative parameter $a$ in \eqref{eq-leaderboard-transform} rescales transformed scores and can affect scale-dependent likelihood and validation quantities, but it cannot alter scale-invariant normality diagnostics. Auxiliary checks across alternative values of $a$ should therefore be interpreted as scale-sensitivity checks, not as evidence for the Gaussian approximation. Additional scale-sensitivity results are available from the authors.

We do not conduct a fully non-Gaussian robustness analysis here, because both the benchmark constructions and the Bayesian estimation procedure rely on the maintained Gaussian structure.
Extending the analysis beyond this setting would require re-specifying the likelihood and re-deriving the benchmark estimators under alternative distributional assumptions.

Overall, the robustness analyses suggest that the two-player reduction is a useful empirical summary of focal rivalry rather than a literal description of full contest dynamics, and that the private-leaderboard validation results should be interpreted as directional evidence under the maintained transformation and Gaussian approximation.


\section{Conclusion}

This paper develops a dynamic game-theoretic model and Bayesian structural estimation framework for online crowdsourcing contests under noisy public feedback. Conditional on focal pairs selected using final private ranks, we recover latent competition states, effort dynamics, and contest-specific primitives from contest-phase submission timing and public leaderboard trajectories, and assess the estimated primitives using private-leaderboard path values excluded from the likelihood. Across 73 Kaggle contests, the diagnostic evidence is strongest for innovation risk; for signal precision, it is directionally consistent but more modest. The counterfactual exercise then traces the model-implied best-output-preserving prize-duration substitution frontier rather than attempting a full welfare comparison of alternative contest designs.


Several limitations deserve mention. First, the empirical setting differs from the baseline model in several respects: many contests use multi-tier or team-based prizes, actual contests involve more than two strategic contestants, leaderboard metrics vary across contests, and the focal pair is selected using final private ranks. Second, the empirical filter assumes that inference variance quickly converges and approximates asynchronous submission-triggered updates with a continuous-update equation. A hybrid filter with discrete measurement updates would represent irregular observation times more faithfully but would substantially reduce tractability. Third, despite the common logit transformation, cross-contest comparisons of innovation risk and private-path volatility may retain a shared transformed-score scale component. These limitations suggest natural extensions involving richer prize structures, multi-agent dynamics, and more flexible learning mechanisms.




% Appendix here
% Options are (1) APPENDIX (with or without general title) or
%             (2) APPENDICES (if it has more than one unrelated sections)
% Outcomment the appropriate case if necessary
%
% \begin{APPENDIX}{<Title of the Appendix>}
% \end{APPENDIX}
%
%   or
%
% \begin{APPENDICES}
% \section{<Title of Section A>}
% \section{<Title of Section B>}
% etc
% \end{APPENDICES}


%\newpage
\begin{APPENDICES}


\section{Proofs}

\proof{Proof of Lemma~\ref{lem-best-output}}
All expectations below are conditional on the fixed initial information $I_0$; we suppress this conditioning in the notation.
Integrating equation~\eqref{eq-individual-output-levels} from $0$ to $T$ gives, for $k\in\{i,j\}$,
\[
x_T^k=\int_0^T m_k(\tilde y_t,t)dt+\frac{\sigma}{\sqrt2}W_T^k.
\]
The equilibrium efforts are integrable by Lemma~\ref{lem-total-effort}, and the Brownian terms have conditional mean zero. Hence,
\[
\mathbb E\left[\frac{x_T^i+x_T^j}{2}\right]
=\frac12\mathbb E\left[\int_0^T
\bigl(m_i(\tilde y_t,t)+m_j(\tilde y_t,t)\bigr)dt\right]
=\frac12M(0,0).
\]
Because $y_T=x_T^i-x_T^j$, the pathwise identity
\[
\max\{x_T^i,x_T^j\}
=\frac{x_T^i+x_T^j}{2}+\frac{|y_T|}{2}
\]
and the preceding expectation establish equation~\eqref{eq-best-output-decomposition}.

We next establish equation~\eqref{eq-terminal-gap-leading}.
Under symmetric costs, $w_i=w_j=w$, and the quantities in Theorem~\ref{thm-equilibrium-effort} reduce to
\[
r_w:=\rho_i=\rho_j=\tanh(w/2),
\qquad
g_w:=\gamma(r_w)=\frac{\sinh w+w}{2}.
\]
The function $\gamma$ is odd, and equation~\eqref{eq-gamma} gives
\[
\gamma'(u)=\frac{2}{(1-u^2)^2}>0,
\qquad u\in(-1,1).
\]
Thus, $\gamma^{-1}$ is odd and strictly increasing. Since
\[
\rho(z)=\gamma^{-1}\!\left(g_w[2\Phi(z)-1]\right)
\]
and $|2\Phi(z)-1|\leq1$, it follows that $|\rho(z)|\leq r_w$ for every $z\in\mathbb R$.
Under symmetry, the definition of $A_+$ in Lemma~\ref{lem-total-effort} therefore implies
\[
\left|A_+(z)e^{-z^2/2}\right|
=\frac{g_w}{\sqrt{2\pi}}
\bigl(1-\rho(z)^2\bigr)|\rho(z)|e^{-z^2/2}
\leq\frac{g_wr_w}{\sqrt{2\pi}}.
\]
Moreover, as $w\downarrow0$,
\[
r_w=\frac{w}{2}+O(w^3),
\qquad
g_w=w+O(w^3),
\]
so $g_wr_w=O(w^2)$.

Let $z_t:=\tilde y_t/(\sigma\sqrt{T-t})$ and define the cumulative equilibrium effort difference
\[
J_T:=\int_0^T
\bigl(m_i(\tilde y_t,t)-m_j(\tilde y_t,t)\bigr)dt.
\]
Subtracting the two effort expressions in equation~\eqref{eq-equilibrium-effort} yields
\[
m_i(\tilde y_t,t)-m_j(\tilde y_t,t)
=\frac{2\sigma}{\sqrt{T-t}}A_+(z_t)e^{-z_t^2/2}.
\]
The preceding uniform bound therefore gives, pathwise,
\[
|J_T|
\leq\frac{2\sigma g_wr_w}{\sqrt{2\pi}}
\int_0^T\frac{dt}{\sqrt{T-t}}
=\frac{4\sigma g_wr_w\sqrt T}{\sqrt{2\pi}}
=O(w^2\sigma\sqrt T).
\]

Finally, $x_0^i=x_0^j=0$ implies $y_0=0$. Integrating the gap dynamics and using $W_T=(W_T^i-W_T^j)/\sqrt2$ gives
\[
y_T=J_T+\sigma W_T,
\qquad W_T\sim\mathcal N(0,T).
\]
Consequently, $\mathbb E|\sigma W_T|=\sigma\sqrt{2T/\pi}$. The pathwise inequality
$\bigl||a+b|-|a|\bigr|\leq|b|$, applied with $a=\sigma W_T$ and $b=J_T$, now yields
\[
\left|\frac12\mathbb E|y_T|
-\frac{\sigma\sqrt T}{\sqrt{2\pi}}\right|
\leq\frac12\mathbb E|J_T|
=O(w^2\sigma\sqrt T),
\]
which proves equation~\eqref{eq-terminal-gap-leading}.
\Halmos
\endproof

\proof{Proof of Proposition~\ref{prop-joint-design}}
For fixed $T>0$, introduce the proof-local functions
\[
a(T):=\frac{\sigma\sqrt T}{\sqrt{2\pi}},
\qquad
b(T):=\frac{\sqrt T}{\sigma c\sqrt{2\pi}},
\qquad
s:=\sigma^2c.
\]
Equation~\eqref{eq-best-output-leading-approximation} can be written as
\[
\bar B(\theta,T)=a(T)+b(T)\theta=b(T)(\theta+s),
\qquad a(T)=s b(T).
\]
Because $V''<0$, $V'$ is strictly decreasing. The limits $V'(0+)=\infty$ and $V'(\infty)=0$ imply that $V'$ maps $(0,\infty)$ onto $(0,\infty)$, so $(V')^{-1}$ is well defined and continuous.

Holding $T$ fixed and differentiating equation~\eqref{eq-designer-best-output-objective} with respect to $\theta$ gives
\[
\frac{\partial\Pi(\theta,T)}{\partial\theta}
=b(T)V'\!\left(a(T)+b(T)\theta\right)-1,
\]
and
\[
\frac{\partial^2\Pi(\theta,T)}{\partial\theta^2}
=b(T)^2V''\!\left(a(T)+b(T)\theta\right)<0.
\]
Thus, the conditional prize problem is strictly concave. Since $V$ is bounded above, $\Pi(\theta,T)\to-\infty$ as $\theta\to\infty$, so its conditional maximizer is unique. It is zero when $b(T)V'(a(T))\leq1$ and otherwise satisfies
\[
V'\!\left(a(T)+b(T)\theta^*(T)\right)=\frac{1}{b(T)}.
\]
The two mutually exclusive cases can be written together as
\[
\begin{aligned}
\theta^*(T)
&=\left[\frac{(V')^{-1}(1/b(T))-a(T)}{b(T)}\right]_+\\
&=\left[\frac{1}{b(T)}(V')^{-1}\!\left(\frac{1}{b(T)}\right)-s\right]_+\\
&=\left[
\sigma c\sqrt{\frac{2\pi}{T}}
(V')^{-1}\!\left(\sigma c\sqrt{\frac{2\pi}{T}}\right)
-\sigma^2c
\right]_+.
\end{aligned}
\]

We next prove existence of a finite positive optimal duration. Define the profile objective
\[
\widehat\Pi(T):=\Pi(\theta^*(T),T),
\qquad T>0.
\]
The explicit expression for $\theta^*(T)$ above implies that $\theta^*(T)$, and hence $\widehat\Pi(T)$, is continuous on $(0,\infty)$.
As $T\downarrow0$, $a(T)\downarrow0$. Concavity of $V$ gives
\[
0\leq a(T)V'(a(T))\leq V(a(T))-V(0)\longrightarrow0.
\]
Since $b(T)=a(T)/s$, we have $b(T)V'(a(T))\to0$, so $\theta^*(T)=0$ for all sufficiently small $T$. Therefore,
\[
\lim_{T\downarrow0}\widehat\Pi(T)=V(0)-D(0).
\]
Moreover, $V'(0+)=\infty$ and continuous differentiability of $D$ at zero imply
\[
\frac{V(a(T))-V(0)}{a(T)}\longrightarrow\infty,
\qquad
\frac{D(T)-D(0)}{a(T)}=O(\sqrt T)\longrightarrow0.
\]
Hence, $\widehat\Pi(T)>V(0)-D(0)$ for all sufficiently small positive $T$, which rules out the zero-duration boundary.

To control the other boundary, choose any $T_1>0$ and set $d_1:=[D(T_1)-D(0)]/T_1>0$. Convexity of $D$ implies, for $T\geq T_1$,
\[
D(T)\geq D(T_1)+d_1(T-T_1),
\]
so $D(T)\to\infty$. Since $V$ is bounded above and $\theta\geq0$,
\[
\widehat\Pi(T)
\leq\sup_{B\geq0}V(B)-D(T)
\longrightarrow-\infty.
\]
The continuity and the two boundary results imply that $\widehat\Pi$ attains a global maximum at some $T^*\in(0,\infty)$. Setting $\theta^*:=\theta^*(T^*)$ proves existence of an optimal pair and, by the conditional-prize formula, equation~\eqref{eq-optimal-prize}.

It remains to establish equation~\eqref{eq-optimal-duration-condition}. Suppose $\theta^*>0$. In a neighborhood of $T^*$, the conditional prize is interior, and its first-order condition is
\[
b(T)V'\!\left(\bar B(\theta^*(T),T)\right)=1.
\]
Because the derivative of the left-hand side with respect to $\theta^*(T)$ is $b(T)^2V''<0$, the implicit function theorem makes $\theta^*(T)$ continuously differentiable in this neighborhood.
When $\theta$ is held fixed, $b'(T)=b(T)/(2T)$ and therefore
\[
\frac{\partial\bar B(\theta,T)}{\partial T}
=b'(T)(\theta+s)
=\frac{\bar B(\theta,T)}{2T}.
\]
Applying the chain rule to the profile objective and then using the prize first-order condition gives the envelope relation
\[
\begin{aligned}
\widehat\Pi'(T)
&=\frac{\partial\Pi}{\partial\theta}
\frac{d\theta^*(T)}{dT}
+\frac{\partial\Pi}{\partial T}\\
&=V'\!\left(\bar B(\theta^*(T),T)\right)
\frac{\bar B(\theta^*(T),T)}{2T}-D'(T).
\end{aligned}
\]
At the interior optimal duration $T^*$, $\widehat\Pi'(T^*)=0$. In addition,
\[
V'\!\left(\bar B(\theta^*,T^*)\right)\bar B(\theta^*,T^*)
=\frac{1}{b(T^*)}b(T^*)(\theta^*+s)
=\theta^*+s.
\]
Substituting this equality into the envelope condition and using $s=\sigma^2c$ gives equation~\eqref{eq-optimal-duration-condition}.
\Halmos
\endproof

\proof{Derivation of Remark~\ref{rem-comparative-statics}}
Retain the proof-local definitions $a(T)=\sigma\sqrt T/\sqrt{2\pi}$, $b(T)=\sqrt T/(\sigma c\sqrt{2\pi})$, and $s=\sigma^2c$. Let
\[
B^*(T):=\bar B(\theta^*(T),T)=b(T)(\theta^*(T)+s).
\]
On the interior-prize region, $b(T)V'(B^*(T))=1$. Since $d\log b(T)/d\log T=1/2$, log-differentiation gives
\[
0=\frac12-R(B^*(T))\frac{d\log B^*(T)}{d\log T},
\]
and hence
\[
\frac{d\log B^*(T)}{d\log T}
=\frac{1}{2R(B^*(T))}.
\]
Using $\theta^*(T)+s=B^*(T)/b(T)$ then yields
\[
\frac{d\log(\theta^*(T)+s)}{d\log T}
=\frac{1-R(B^*(T))}{2R(B^*(T))}\geq0,
\]
where the inequality follows from $R\leq1$. At a prize corner, the marginal return to increasing the prize from zero is $b(T)V'(a(T))$. Its log derivative is
\[
\frac{d\log[b(T)V'(a(T))]}{d\log T}
=\frac{1-R(a(T))}{2}\geq0.
\]
Thus, the set on which the prize is interior is an upper interval in $T$, and continuity at its boundary implies that $\theta^*(T)$ is nondecreasing on $(0,\infty)$.

Now define the gross-value profile
\[
G(T):=\max_{\theta\geq0}
\{V(\bar B(\theta,T))-\theta\}.
\]
Within either prize region, the envelope theorem gives
\[
G'(T)=V'(B^*(T))\frac{B^*(T)}{2T}.
\]
At a prize corner, $B^*(T)=a(T)$, so
\[
\frac{d\log G'(T)}{d\log T}
=-1+\frac{1-R(a(T))}{2}
=-\frac{1+R(a(T))}{2}<0.
\]
On the interior-prize region, the prize first-order condition gives
\[
G'(T)=\frac{\theta^*(T)+s}{2T},
\]
and therefore
\[
\frac{d\log G'(T)}{d\log T}
=\frac{1-R(B^*(T))}{2R(B^*(T))}-1
=\frac{1-3R(B^*(T))}{2R(B^*(T))}<0,
\]
where strict negativity follows from $R>1/3$. At a boundary between the two prize regions, $\theta^*(T)=0$ and $b(T)V'(a(T))=1$, so both expressions for $G'(T)$ equal $s/(2T)$. Hence, $G'(T)$ is continuous and strictly decreasing. Since convexity of $D$ makes $D'(T)$ nondecreasing, the derivative $G'(T)-D'(T)$ crosses zero at most once. Existence from Proposition~\ref{prop-joint-design} then implies that the optimal duration is unique.
\Halmos
\endproof

\proof{Proof of Lemma~\ref{lmm-params-identifiability}}
The proof proceeds under the maintained structural specification in
(\ref{eq-equilibrium-effort}), (\ref{eq-filtered-y-update}),
(\ref{eq-ihpp-prob}), and (\ref{eq-obs_gap-prob}).
Suppose two parameter vectors $\Theta=(c_i,c_j,\mu_0,\sigma,\lambda)$ and
$\Theta'=(c_i',c_j',\mu_0',\sigma',\lambda')$ induce the same likelihood in a
neighborhood of the true parameter value.

First, consider the Gaussian leaderboard component in (\ref{eq-obs_gap-prob}). 
Equality of Gaussian likelihoods for all sampled paths implies equality of both covariance and mean structures. 
Hence,
\begin{equation*}
\Sigma_y + D_\lambda
=
\Sigma_y' + D_{\lambda'},
\qquad
\mu_k=\mu_k'\ \text{for all } k.
\end{equation*}
Here $D_\lambda$ and $D_{\lambda'}$ are diagonal matrices with $D_{\lambda,kk}=1/(\lambda h_k)$ and $D_{\lambda',kk}=1/(\lambda' h_k)$, where the readout intervals $\{h_k\}_{k=1}^N$ are fixed by the observation design.
With at least two distinct leaderboard observations, the off-diagonal covariance terms equal $\sigma^2\min(t_i,t_j)$ and are unaffected by the independent signal-noise term. 
Therefore equality of the off-diagonal covariance terms implies $\sigma=\sigma'$. Given $\sigma=\sigma'$, equality of the diagonal terms then implies $\lambda=\lambda'$.

It remains to identify $(c_i,c_j,\mu_0)$.
Given $(\sigma,\lambda)$ and fixed $r$, define $\vartheta=(c_i,c_j,\mu_0)$.
For any candidate $\vartheta$, the filtering recursion and equilibrium effort rule generate a perceived-state path $\tilde y_t(\vartheta)$, effort paths $m_i(t;\vartheta)$ and $m_j(t;\vartheta)$, submission intensities
\[
\tau_i(t;\vartheta)=r\,m_i(t;\vartheta),
\qquad
\tau_j(t;\vartheta)=r\,m_j(t;\vartheta),
\]
and public-gap mean values $\mu_k(\vartheta)$ at the observed leaderboard update times.
Let
\[
G(\vartheta)
=
\left(
\{\tau_i(t;\vartheta)\}_{t\in\mathcal T_i},
\{\tau_j(t;\vartheta)\}_{t\in\mathcal T_j},
\{\mu_k(\vartheta)\}_{k=1}^N
\right)
\]
denote this reduced-form prediction map, evaluated on the support of observed submission and update times.

We impose the full-rank variation condition that the Jacobian
$\partial G(\vartheta_0)/\partial\vartheta$ has full column rank at the true parameter value $\vartheta_0$.
This condition formalizes the requirement that the realized public-history support contain enough variation to distinguish changes in effort costs from changes in the initial perceived gap.
Under this condition, equality of the point-process intensity functions and equality of the Gaussian mean structure imply $G(\vartheta)=G(\vartheta_0)$ only locally at $\vartheta=\vartheta_0$.
Hence $(c_i,c_j,\mu_0)$ is locally identified.

Combining this local identification of $(c_i,c_j,\mu_0)$ with the covariance-based identification of $(\sigma,\lambda)$ yields local identification of $(c_i,c_j,\sigma,\lambda,\mu_0)$ conditional on fixed $r$.
\Halmos
\endproof






\section{Auxiliary Results}

\subsection{Solve $S_t$ in Equation (\ref{filtered-S})}\label{app-S-equ}

If $S_t$ is in steady state, $dS_t/dt = 0$ $\Leftrightarrow$ $S_t = \bar{S} \equiv \sigma/\sqrt{\lambda}$.
If $S_t$ is not in steady state, i.e., $S_t\neq\bar{S}$, we first isolate the two variables and get
\begin{equation*}
	\frac{dS}{\sigma^2-\lambda S^2} = dt
\end{equation*}
Integrating both sides gives
\begin{equation*}
    t = \int \frac{dS}{\sigma^2-\lambda S^2}
    = \frac{1}{\sigma\sqrt{\lambda}}
    \int \frac{d(S\sqrt{\lambda}/\sigma)}{1-(S\sqrt{\lambda}/\sigma)^2}
    \equiv \frac{1}{\sigma\sqrt{\lambda}}\int \frac{du}{1-u^2},
\end{equation*}
where $u = S\sqrt{\lambda}/\sigma = S/\bar{S}$.
For $\tanh^{-1}(u)$ we require $|u|<1$ ($S<\bar S$), and for $\coth^{-1}(u)$ we require $|u|>1$ ($S>\bar S$).
Hence,
\begin{equation*}
    \sigma\sqrt{\lambda}\cdot t = \begin{cases}
        \tanh^{-1}(u)-K_1, & \text{if } |u|<1\\
        \coth^{-1}(u)-K_2, & \text{if } |u|>1
    \end{cases} = \begin{cases}
        \tanh^{-1}(S/\bar{S})-K_1, & \text{if } S<\bar{S}\\
        \coth^{-1}(S/\bar{S})-K_2, & \text{if } S>\bar{S}
\end{cases}
\end{equation*}
Thus, in the non-steady-state case,
\begin{equation*}
	S = \begin{cases}
		\bar{S}\cdot\tanh(\sigma\sqrt{\lambda}\cdot t + K_1), & \text{if } S<\bar{S}\\
		\bar{S}\cdot\coth(\sigma\sqrt{\lambda}\cdot t + K_2), & \text{if } S>\bar{S}
	\end{cases}
\end{equation*}
Finally, we determine the constants $K_1$, $K_2$ by the initial condition $S_0$ and have
\begin{align*}
	K_1 = \tanh^{-1}(S_0/\bar{S}),\quad
	K_2 = \coth^{-1}(S_0/\bar{S})
\end{align*}
If $\lambda = 0$, we have $S_t = S_0+\sigma^2t$, i.e., the estimation variance increases linearly over time.
If $\lambda > 0$, the solution of (\ref{filtered-S}) is
\begin{equation*}
	S_t =
	\begin{cases}
        \bar{S} \cdot \tanh\left\{t\cdot\sigma\sqrt{\lambda} +  \tanh^{-1}\left(S_0\big/\bar{S}\right)\right\} & \text{if } S_0 < \bar{S} \\
        \bar{S} & \text{if } S_0 = \bar{S} \\
        \bar{S} \cdot \coth\left\{t\cdot\sigma\sqrt{\lambda} +  \coth^{-1}\left(S_0\big/\bar{S}\right)\right\} & \text{if } S_0 > \bar{S}
	\end{cases}
\end{equation*}
Specifically, $\bar{S} = \sigma / \sqrt{\lambda}$ when $\lambda>0$ and $\bar{S} = \infty$ when $\lambda = 0$.


\subsection{Consistency Benchmark for $\sigma^2$}\label{appexidx-normal-benchmark}

This appendix derives the private-leaderboard benchmark for $\sigma$ used in Section~\ref{sect-validation}.
For a given contest, let $\Delta y_k=y_k-y_{k-1}$ denote the private-leaderboard gap increment between two adjacent update times, and let $h_k=t_k-t_{k-1}$ be the corresponding elapsed time.
The posterior effort trajectories provide a plug-in drift adjustment
\[
\hat\mu_k=\int_{t_{k-1}}^{t_k}\left[\hat m_i(s)-\hat m_j(s)\right]ds.
\]
Under the diffusion approximation, the drift-adjusted increment satisfies
\[
\Delta y_k-\hat\mu_k\sim\mathcal N(0,\sigma^2 h_k),
\qquad k=1,\ldots,n.
\]
Treating these increments as approximately independent, the log-likelihood for $\sigma^2$ is, up to constants,
\[
\ell(\sigma^2)
=
-\frac{n}{2}\log(\sigma^2)
-\frac{1}{2}\sum_{k=1}^n\log h_k
-\frac{1}{2\sigma^2}\sum_{k=1}^n
\frac{(\Delta y_k-\hat\mu_k)^2}{h_k}.
\]
Differentiating with respect to $\sigma^2$ and setting the derivative equal to zero gives
\[
-\frac{n}{2\sigma^2}
+
\frac{1}{2(\sigma^2)^2}
\sum_{k=1}^n
\frac{(\Delta y_k-\hat\mu_k)^2}{h_k}
=0.
\]
Therefore,
\begin{equation}\label{normal-benchmark-sigma}
\hat{\sigma}^2
=
\frac{1}{n}\sum_{k=1}^n
\frac{(\Delta y_k-\hat\mu_k)^2}{h_k}.
\end{equation}
In the empirical application, $\hat\mu_k$ is computed from the Bayesian posterior effort paths, so this benchmark is a plug-in private-leaderboard consistency check rather than a fully model-free estimator. Moreover, asynchronous update timing means that adjacent displayed gaps may combine a new score for one contestant with a carried-forward score for the other. The expression in \eqref{normal-benchmark-sigma} is therefore an approximate benchmark for the carried-forward gap process, not an exact maximum likelihood estimator for synchronously observed latent diffusion increments.








\section{Statistics}\label{appendix-statistics}

Table~\ref{tab-appendix-full-results} reports raw contest-level estimates for transparency, even though some benchmark values are highly dispersed.
Large values of $\lambda^{\text{bench}}$ can arise mechanically when public-private score discrepancies are very small, because the benchmark is inverse-quadratic in those discrepancies.
Accordingly, cross-contest comparisons of $\lambda^{\text{bench}}$ are most interpretable on a transformed scale (as used in the main-text regressions).

\begin{longtable}{cccccccccccrc}
\caption{Contest-Level Structural Estimates and Benchmarks}\label{tab-appendix-full-results}\\
\toprule
\multirow{2}{*}{\makecell{\rule{0pt}{2.5ex}Contest\\Id}} &
\multirow{2}{*}{\makecell{\rule{0pt}{2.5ex}$N_i$}} &
\multirow{2}{*}{\makecell{\rule{0pt}{2.5ex}$N_j$}} &
\multirow{2}{*}{\makecell{\rule{0pt}{2.5ex}Prize\\(k, USD)}} &
\multirow{2}{*}{\makecell{\rule{0pt}{2.5ex}Public\\Data \%}} &
\multicolumn{6}{c}{Bayesian Posterior Mean} &
\multicolumn{2}{c}{Diagnostic Benchmarks} \\
\cmidrule(lr){6-11}
\cmidrule(lr){12-13}
 &  &  &  &  & $\lambda$ & $\sigma$ & $\mu_0$ & $c_i$ & $c_j$ & $r$ & \makecell{$\lambda^\text{bench}$} & $\sigma^\text{bench}$ \\
\midrule
2445 & 7 & 19 & 5.00 & 25 & 15.377 & 0.980 & -0.089 & 4.304 & 3.154 & 2.088 & 2064.402 & 0.394 \\
2454 & 19 & 13 & 0.15 & 62 & 15.016 & 1.030 & 4.377 & 0.688 & 0.810 & 18.035 & 768.208 & 0.659 \\
2464 & 11 & 12 & 0.95 & 20 & 9.844 & 2.325 & -0.474 & 1.631 & 1.380 & 14.217 & 95.868 & 3.741 \\
2478 & 13 & 11 & 0.95 & 30 & 13.636 & 0.673 & -0.444 & 3.171 & 3.811 & 6.333 & 12004.117 & 0.770 \\
2489 & 12 & 21 & 0.50 & 10 & 16.650 & 0.576 & 0.431 & 3.535 & 2.303 & 11.210 & 5663.716 & 0.843 \\
2549 & 13 & 15 & 0.50 & 30 & 12.214 & 1.357 & -1.055 & 1.409 & 1.157 & 14.758 & 520.220 & 2.891 \\
2551 & 48 & 30 & 1.50 & 30 & 26.143 & 0.539 & -0.286 & 3.207 & 4.141 & 6.669 & 93276.660 & 0.512 \\
2667 & 55 & 67 & 30.00 & 50 & 16.977 & 1.695 & 0.758 & 3.919 & 3.803 & 1.184 & 3.650 & 5.430 \\
2749 & 22 & 22 & 3.00 & 42 & 20.708 & 0.701 & 0.906 & 4.308 & 4.022 & 4.367 & 42318.566 & 0.340 \\
2762 & 12 & 21 & 1.00 & 35 & 11.659 & 0.779 & 1.068 & 3.103 & 2.237 & 13.005 & 45.399 & 1.244 \\
2860 & 9 & 6 & 1.00 & 30 & 12.735 & 0.816 & 0.133 & 3.096 & 3.728 & 6.933 & 22731.842 & 0.034 \\
3064 & 13 & 5 & 0.20 & 25 & 14.106 & 0.630 & -0.788 & 1.160 & 2.780 & 13.821 & 8157921.303 & 0.263 \\
3080 & 25 & 27 & 0.09 & 25 & 23.442 & 0.531 & -1.217 & 0.623 & 0.643 & 19.327 & 149115.198 & 0.187 \\
3288 & 20 & 10 & 1.00 & 0 & 10.864 & 3.729 & -1.169 & 0.822 & 1.750 & 15.448 & 3.518 & 0.616 \\
3338 & 31 & 29 & 1.50 & 30 & 24.827 & 0.541 & 0.203 & 3.508 & 3.345 & 7.687 & 9639.132 & 0.244 \\
3353 & 26 & 17 & 10.00 & 30 & 16.307 & 1.232 & -0.249 & 3.749 & 4.305 & 2.066 & 23439.960 & 1.410 \\
3366 & 16 & 11 & 0.50 & 50 & 15.653 & 0.693 & -1.936 & 2.055 & 3.018 & 13.340 & 4647.231 & 0.499 \\
3507 & 9 & 13 & 0.50 & 33 & 14.030 & 1.055 & 0.529 & 2.440 & 1.891 & 14.095 & 159.184 & 1.101 \\
3509 & 8 & 13 & 0.20 & 30 & 12.714 & 0.741 & -0.024 & 2.491 & 1.513 & 13.807 & 1802.071 & 0.884 \\
3526 & 15 & 13 & 5.00 & 20 & 15.809 & 0.793 & -0.053 & 4.127 & 4.159 & 1.280 & 6874.479 & 0.242 \\
3641 & 7 & 41 & 1.50 & 43 & 20.406 & 0.579 & 0.615 & 4.292 & 1.875 & 3.434 & 2372.810 & 0.319 \\
3774 & 29 & 26 & 0.50 & 50 & 23.554 & 0.537 & -0.997 & 2.580 & 2.898 & 13.300 & 38.059 & 0.952 \\
3800 & 21 & 28 & 3.00 & 15 & 20.046 & 0.820 & 3.384 & 4.322 & 3.599 & 6.632 & 1094.443 & 0.625 \\
3926 & 25 & 20 & 0.30 & 25 & 22.016 & 0.547 & -0.047 & 1.237 & 1.529 & 13.393 & 4272.601 & 0.332 \\
3928 & 15 & 21 & 0.68 & 70 & 19.525 & 0.584 & -0.421 & 3.604 & 3.051 & 11.838 & 81382.464 & 0.139 \\
3929 & 35 & 31 & 8.00 & 50 & 5.865 & 3.653 & 0.062 & 3.165 & 3.520 & 10.320 & 4021.718 & 5.862 \\
3960 & 47 & 50 & 8.00 & 40 & 28.973 & 0.883 & -0.992 & 4.624 & 4.670 & 3.583 & 589.028 & 0.514 \\
%4031 & 6 & 61 & 5.00 & 30 & 19.132 & 0.831 & 0.605 & 4.487 & 1.975 & 3.021 & 125709.402 & 1.214 \\
4104 & 8 & 6 & 20.00 & 20 & 10.864 & 2.116 & 0.036 & 3.729 & 3.994 & 0.742 & 36704.448 & 0.141 \\
4366 & 31 & 31 & 8.00 & 2 & 23.414 & 0.913 & -0.434 & 4.367 & 4.554 & 2.979 & 3592.689 & 0.487 \\
4407 & 44 & 51 & 4.00 & 19 & 14.573 & 1.670 & 0.070 & 3.591 & 3.238 & 12.165 & 103.661 & 2.693 \\
4453 & 15 & 18 & 2.00 & 30 & 17.165 & 0.617 & 1.204 & 4.110 & 3.071 & 2.607 & 44736.399 & 0.892 \\
4477 & 7 & 11 & 2.00 & 50 & 11.870 & 1.104 & -3.305 & 3.867 & 3.222 & 5.142 & 7249.061 & 0.325 \\
4488 & 11 & 8 & 2.00 & 1 & 8.033 & 5.347 & 1.108 & 2.160 & 2.843 & 12.739 & 1301.149 & 1.559 \\
4493 & 19 & 14 & 10.00 & 40 & 16.683 & 1.161 & -0.191 & 3.883 & 4.306 & 1.567 & 7848.309 & 0.355 \\
4657 & 66 & 62 & 4.00 & 30 & 34.581 & 0.559 & 0.330 & 4.657 & 4.691 & 3.410 & 1143492.186 & 0.205 \\
4699 & 13 & 22 & 5.00 & 30 & 18.734 & 0.765 & 0.112 & 4.242 & 3.833 & 1.207 & 154786.619 & 0.211 \\
4986 & 16 & 5 & 10.00 & 50 & 13.378 & 1.203 & -0.029 & 3.287 & 4.306 & 1.214 & 66271.436 & 0.393 \\
5056 & 18 & 35 & 5.00 & 33 & 21.957 & 0.720 & -0.780 & 4.484 & 4.094 & 2.763 & 564328.562 & 0.324 \\
5144 & 11 & 9 & 20.00 & 20 & 7.244 & 3.118 & 10.283 & 3.617 & 3.737 & 1.411 & 3589.357 & 2.695 \\
5174 & 42 & 47 & 3.00 & 50 & 23.823 & 0.744 & -1.240 & 4.346 & 4.162 & 8.369 & 233854.652 & 1.622 \\
5229 & 28 & 31 & 15.00 & 10 & 12.332 & 3.417 & 0.320 & 3.961 & 3.564 & 6.028 & 12.149 & 2.880 \\
5261 & 45 & 43 & 10.00 & 30 & 26.183 & 0.965 & -1.578 & 4.361 & 4.651 & 2.525 & 81980.068 & 0.965 \\
5357 & 36 & 22 & 5.00 & 30 & 22.768 & 0.774 & 1.093 & 4.028 & 4.506 & 3.022 & 12965.328 & 0.354 \\
5390 & 22 & 17 & 4.00 & 30 & 17.724 & 0.819 & 0.513 & 3.650 & 4.317 & 4.262 & 1225.164 & 1.086 \\
5497 & 44 & 24 & 4.00 & 30 & 25.292 & 0.652 & 0.642 & 4.231 & 4.571 & 3.005 & 11114241.101 & 0.233 \\
6322 & 23 & 26 & 10.00 & 66 & 9.106 & 2.653 & -0.405 & 3.753 & 3.536 & 5.852 & 59.531 & 3.981 \\
6927 & 19 & 16 & 4.00 & 1 & 18.477 & 0.918 & -0.601 & 3.925 & 4.247 & 4.178 & 8619.418 & 0.272 \\
7115 & 19 & 39 & 4.00 & 30 & 15.827 & 0.848 & 0.444 & 4.314 & 3.226 & 3.401 & 314182.624 & 1.469 \\
7162 & 51 & 14 & 1.00 & 50 & 20.424 & 0.808 & -1.269 & 1.209 & 4.010 & 8.932 & 42025.044 & 1.093 \\
7634 & 52 & 44 & 2.00 & 2 & 23.259 & 0.568 & 0.761 & 3.767 & 4.307 & 8.242 & 7551.190 & 1.509 \\
7878 & 5 & 7 & 4.00 & 29 & 9.554 & 1.484 & -0.612 & 3.815 & 3.285 & 1.450 & 360.702 & 0.875 \\
8076 & 34 & 29 & 6.00 & 10 & 22.880 & 0.806 & -1.387 & 4.321 & 4.419 & 2.903 & 10417.833 & 1.134 \\
8078 & 29 & 32 & 4.00 & 36 & 8.066 & 1.796 & 1.664 & 3.587 & 3.088 & 10.469 & 112.042 & 5.352 \\
8219 & 16 & 23 & 0.40 & 22 & 15.559 & 0.604 & -0.087 & 2.601 & 1.940 & 14.337 & 726441.721 & 1.542 \\
8396 & 53 & 59 & 0.50 & 1 & 13.606 & 1.385 & 7.016 & 0.880 & 0.767 & 16.517 & 1425.343 & 2.451 \\
%8540 & 30 & 38 & 5.00 & 18 & 22.939 & 0.840 & -0.117 & 4.326 & 4.212 & 4.251 & 7531.700 & 1.251 \\
9120 & 52 & 44 & 10.00 & 20 & 28.586 & 0.876 & -0.006 & 4.690 & 4.746 & 1.789 & 47120.478 & 0.219 \\
9949 & 21 & 35 & 5.00 & 20 & 21.022 & 0.923 & 1.016 & 4.539 & 3.565 & 3.200 & 14751.471 & 0.563 \\
10200 & 22 & 25 & 4.00 & 9 & 15.905 & 1.513 & -7.127 & 3.619 & 3.583 & 6.466 & 34488.640 & 1.034 \\
10684 & 16 & 58 & 4.00 & 57 & 18.929 & 0.659 & -7.058 & 4.450 & 3.642 & 4.785 & 35.867 & 1.342 \\
13333 & 18 & 19 & 2.00 & 25 & 19.948 & 0.638 & 0.168 & 3.975 & 3.913 & 6.246 & 77062.947 & 0.079 \\
14242 & 68 & 37 & 3.00 & 20 & 31.135 & 0.558 & 0.383 & 4.266 & 4.663 & 5.096 & 2974.188 & 0.372 \\
14420 & 56 & 20 & 20.00 & 22 & 13.393 & 1.571 & -3.492 & 3.144 & 4.609 & 1.464 & 5.221 & 5.899 \\
18045 & 40 & 50 & 4.00 & 30 & 26.932 & 0.654 & -3.178 & 4.269 & 4.344 & 3.532 & 101.449 & 0.809 \\
19018 & 63 & 67 & 8.00 & 30 & 29.425 & 0.835 & 1.891 & 4.634 & 4.360 & 2.918 & 18132.000 & 0.863 \\
19991 & 22 & 50 & 4.00 & 20 & 24.460 & 0.658 & 0.062 & 4.599 & 3.990 & 3.472 & 934.865 & 1.058 \\
20270 & 20 & 16 & 4.00 & 30 & 18.848 & 0.815 & -0.255 & 3.707 & 4.250 & 3.938 & 285.531 & 0.531 \\
21669 & 38 & 37 & 1.00 & 21 & 17.621 & 0.831 & -1.649 & 1.917 & 2.520 & 14.959 & 2068.038 & 2.604 \\
22962 & 23 & 18 & 4.00 & 24 & 17.168 & 0.879 & -1.295 & 3.760 & 4.158 & 4.583 & 7451.091 & 1.143 \\
23249 & 21 & 40 & 1.00 & 16 & 24.123 & 0.532 & 0.538 & 2.978 & 2.075 & 7.226 & 1758578.906 & 0.101 \\
23652 & 24 & 24 & 1.00 & 20 & 20.607 & 0.643 & -0.193 & 2.862 & 3.456 & 14.124 & 76202.298 & 1.056 \\
37077 & 11 & 46 & 4.00 & 24 & 19.329 & 0.886 & 0.584 & 4.311 & 2.254 & 3.507 & 8747.907 & 1.563 \\
38128 & 27 & 52 & 8.00 & 42 & 25.657 & 0.808 & -0.562 & 4.687 & 4.444 & 1.655 & 25574.217 & 0.263 \\
38760 & 10 & 17 & 5.00 & 33 & 12.412 & 1.155 & 0.026 & 4.192 & 3.255 & 2.936 & 5717587.740 & 2.551 \\
\bottomrule
\end{longtable}

\end{APPENDICES}



% Acknowledgments here
\ACKNOWLEDGMENT{The authors gratefully acknowledge the valuable research assistance provided by Simian Zhang during the course of this study.
}


% References here (outcomment the appropriate case)

% CASE 1: BiBTeX used to constantly update the references
%   (while the paper is being written).
%\bibliographystyle{informs2014} % outcomment this and next line in Case 1
%\bibliography{<your bib file(s)>} % if more than one, comma separated

% CASE 2: BiBTeX used to generate mypaper.bbl (to be further fine tuned)
%\input{mypaper.bbl} % outcomment this line in Case 2

%If you don't use BiBTex, you can manually itemize references as shown below.


%\bibliographystyle{nonumber}
\bibliographystyle{informs2014}
\bibliography{_Literatures.bib}

%%%%%%%%%%%%%%%%%
\end{document}
%%%%%%%%%%%%%%%%%
